\chapter{Microphysics and Elastodynamics of the FCC Active Vacuum}
\label{app:fcc_microphysics}

\noindent\textbf{Subject:} Microscopic derivation of the FCC contact action, endpoint transfer closure, dynamic matrix homogenization, and interstitial condensate response.

\section{Constrained Transverse FCC Contact Action and Local Dispersion}

\subsection{Scope and Constitutive Setting}

This section provides a minimal conservative action for transverse state transfer through the active FCC Soft Floor under a locally quasistatic gravitational strain. Its purpose is to determine the dynamical consequences of the endpoint-transfer closure developed in Section 10.4.

The analysis is restricted to the active FCC phase,
\begin{equation}
    0 \leq \alpha_{\mathrm{G}} < 1,
\end{equation}
a gravitational background that varies negligibly during one lattice oscillation,
\begin{equation}
    \dot{\alpha}_{\mathrm{G}} \approx 0,
\end{equation}
and a coherent transverse mode for which the local phase-mismatch factor satisfies
\begin{equation}
    \cos(\Delta\varphi) \approx 1.
\end{equation}

The solid FCC scaffold supplies the transverse restorative tension, the interstitial condensate supplies the dynamic surface inertia, and the phase-locked Hard Floor supplies the vertical constraint represented macroscopically by $\alpha_{\mathrm{G}}$.

The derivation does not independently determine the endpoint-transfer closure. Rather, it establishes the conservative contact dynamics and dispersion relation that follow when that closure is imposed as the local kinematic constraint.

\subsection{Nearest-Neighbor Transverse Contact}

For touching nodes in the FCC Soft Floor, the nearest-neighbor separation is
\begin{equation}
    d_{\mathrm{nn}} = 2r_{\mathrm{D}}.
\end{equation}

Consider a selected nearest-neighbor propagation path and one transverse polarization orthogonal to the bond direction.

Let $u_n(t)$ denote the completed transverse displacement of the $n$-th Soft Floor node. The transverse strain carried by the contact between nodes $n$ and $n+1$ is
\begin{equation}
    \varepsilon_{n,n+1}^{\perp} = \frac{u_{n+1}-u_n}{d_{\mathrm{nn}}}.
\end{equation}

Because $\tau_{\mathrm{D}}$ is the baseline surface shear tension of the active Soft Floor, the incremental contact energy per unit participating interface area is
\begin{equation}
    \boxed{ \mathcal{U}_{n,n+1}^{\perp} = \frac{\tau_{\mathrm{D}}}{2d_{\mathrm{nn}}^2} \left( u_{n+1}-u_n \right)^2 }.
\end{equation}

Its two-node potential-energy Hessian is
\begin{equation}
    \boxed{ \mathbf{H}_{U}^{(n,n+1)} = \frac{\tau_{\mathrm{D}}}{d_{\mathrm{nn}}^2}
    \begin{pmatrix}
        1 & -1 \\
        -1 & 1
    \end{pmatrix} }.
\end{equation}

At the retained constitutive order, gravitational loading does not modify this transverse potential Hessian:
\begin{equation}
    \boxed{ \mathbf{H}_{U}^{(n,n+1)}(\alpha_{\mathrm{G}}) = \mathbf{H}_{U}^{(n,n+1)}(0) }.
\end{equation}

The strain dependence therefore enters through the constrained kinetic sector rather than through continuous softening of the baseline FCC shear tension.

\subsection{Latent and Completed Reconfiguration Coordinates}

Let $y_n(t)$ denote the condensate-supported reconfiguration amplitude executed locally by the $n$-th node. This coordinate represents the underlying node--condensate motion required to produce a completed transverse displacement $u_n$.

In the unstrained Soft Floor,
\begin{equation}
    \alpha_{\mathrm{G}} = 0,
\end{equation}
the full reconfiguration is transmitted, and therefore
\begin{equation}
    u_n = y_n.
\end{equation}

For a nearest-neighbor transfer from node $i$ to node $j$, Section 10.4 defines the endpoint-transfer factor
\begin{equation}
    \boxed{ \mathcal{C}_{ij} = \left[ 1-\alpha_{\mathrm{G}}^{(i)} \right] \left[ 1-\alpha_{\mathrm{G}}^{(j)} \right] }.
\end{equation}

Within a locally uniform region,
\begin{equation}
    \alpha_{\mathrm{G}}^{(i)} \approx \alpha_{\mathrm{G}}^{(j)} \approx \alpha_{\mathrm{G}},
\end{equation}
and hence
\begin{equation}
    \mathcal{C}_{ij} \approx \left( 1-\alpha_{\mathrm{G}} \right)^2.
\end{equation}

For the locally uniform directional reduction considered below, write this common factor as
\begin{equation}
    \mathcal{C} = \left( 1-\alpha_{\mathrm{G}} \right)^2.
\end{equation}

The completed displacement is constrained to the underlying reconfiguration amplitude by
\begin{equation}
    \boxed{ u_n = \mathcal{C} y_n }.
\end{equation}

Equivalently,
\begin{equation}
    y_n = \frac{u_n}{\mathcal{C}}.
\end{equation}

The transfer factor $\mathcal{C}$ is therefore the local Jacobian between the underlying condensate-supported reconfiguration and the completed displacement of the active FCC state.

\subsection{Constrained Transverse Action}

The baseline surface inertia of the Soft Floor is $\mu_{\mathrm{D}}$. The kinetic energy per unit participating interface area associated with the underlying reconfiguration coordinate is therefore
\begin{equation}
    \mathcal{T}_n = \frac{1}{2} \mu_{\mathrm{D}}\dot{y}_n^2.
\end{equation}

The transverse Lagrangian per unit participating interface area is
\begin{equation}
    \boxed{ \mathscr{L}_{\perp} = \sum_n \left[ \frac{1}{2}\mu_{\mathrm{D}}\dot{y}_n^2 - \frac{\tau_{\mathrm{D}}}{2d_{\mathrm{nn}}^2} \left( u_{n+1}-u_n \right)^2 \right] }
\end{equation}
subject to the holonomic constraint
\begin{equation}
    u_n - \mathcal{C} y_n = 0.
\end{equation}

The Hard Floor does not appear as an ordinary finite-frequency oscillator in this action. Its effective rigidity and vanishing intrinsic inertia place it in the singular constraint limit. The Tension Column therefore restricts the admissible relation between $y_n$ and $u_n$ without introducing an independent propagating substrate degree of freedom.

For a quasistatic gravitational background,
\begin{equation}
    \dot{\mathcal{C}} = 0,
\end{equation}
and the constraint gives
\begin{equation}
    \dot{y}_n = \frac{\dot{u}_n}{\mathcal{C}}.
\end{equation}

Eliminating $y_n$ yields the reduced transverse Lagrangian
\begin{equation}
    \boxed{ \mathscr{L}_{\perp}^{(\mathrm{red})} = \sum_n \left[ \frac{1}{2} \frac{\mu_{\mathrm{D}}}{\mathcal{C}^2} \dot{u}_n^2 - \frac{\tau_{\mathrm{D}}}{2d_{\mathrm{nn}}^2} \left( u_{n+1}-u_n \right)^2 \right] }.
\end{equation}

The reduced kinetic Hessian is consequently
\begin{equation}
    \boxed{ \mathbf{H}_T(\alpha_{\mathrm{G}}) = \frac{\mu_{\mathrm{D}}}{\mathcal{C}^2}\mathbf{I} }.
\end{equation}

Thus, the effective surface inertia is
\begin{equation}
    \boxed{ \mu_{\mathrm{eff}} = \frac{\mu_{\mathrm{D}}}{\mathcal{C}^2} }.
\end{equation}

Using
\begin{equation}
    \mathcal{C} = \left( 1-\alpha_{\mathrm{G}} \right)^2,
\end{equation}
one obtains
\begin{equation}
    \boxed{ \mu_{\mathrm{eff}} = \frac{\mu_{\mathrm{D}}}{ \left( 1-\alpha_{\mathrm{G}} \right)^4 } }.
\end{equation}

The inverse-fourth-power law is therefore the reduced kinetic consequence of the two-endpoint transfer Jacobian.

\subsection{Discrete Equation of Motion}

Variation of the reduced action with respect to $u_n$ gives
\begin{equation}
    \frac{\mu_{\mathrm{D}}}{\mathcal{C}^2}\ddot{u}_n = \frac{\tau_{\mathrm{D}}}{d_{\mathrm{nn}}^2} \left( u_{n+1} - 2u_n + u_{n-1} \right).
\end{equation}

Equivalently,
\begin{equation}
    \boxed{ \mu_{\mathrm{eff}}\ddot{u}_n = \frac{\tau_{\mathrm{D}}}{d_{\mathrm{nn}}^2} \left( u_{n+1} - 2u_n + u_{n-1} \right) }.
\end{equation}

The discrete Laplacian on the right-hand side is the nearest-neighbor transverse restoring force of the FCC contact sequence.

Consider a plane-wave mode
\begin{equation}
    u_n(t) = u_0 \exp \left[ i \left( knd_{\mathrm{nn}} - \omega t \right) \right].
\end{equation}

The nearest-neighbor difference becomes
\begin{equation}
    u_{n+1}-2u_n+u_{n-1} = -4u_n \sin^2 \left( \frac{kd_{\mathrm{nn}}}{2} \right).
\end{equation}

Substitution gives the dispersion relation
\begin{equation}
    \boxed{ \omega^2(k,\alpha_{\mathrm{G}}) = \frac{4\tau_{\mathrm{D}}}{\mu_{\mathrm{D}}d_{\mathrm{nn}}^2} \mathcal{C}^2 \sin^2 \left( \frac{kd_{\mathrm{nn}}}{2} \right) }.
\end{equation}

Using the baseline Soft Floor identity
\begin{equation}
    c_{\mathrm{E}}^2 = \frac{\tau_{\mathrm{D}}}{\mu_{\mathrm{D}}},
\end{equation}
this becomes
\begin{equation}
    \boxed{ \omega^2(k,\alpha_{\mathrm{G}}) = \frac{4c_{\mathrm{E}}^2\mathcal{C}^2}{d_{\mathrm{nn}}^2} \sin^2 \left( \frac{kd_{\mathrm{nn}}}{2} \right) }.
\end{equation}

For a locally uniform gravitational strain,
\begin{equation}
    \mathcal{C} = \left( 1-\alpha_{\mathrm{G}} \right)^2,
\end{equation}
and therefore
\begin{equation}
    \boxed{ \omega^2(k,\alpha_{\mathrm{G}}) = \frac{4c_{\mathrm{E}}^2}{d_{\mathrm{nn}}^2} \left( 1-\alpha_{\mathrm{G}} \right)^4 \sin^2 \left( \frac{kd_{\mathrm{nn}}}{2} \right) }.
\end{equation}

This retains the ordinary nearest-neighbor lattice dispersion while rescaling its characteristic frequency through the gravitational endpoint constraint.

\subsection{Long-Wavelength Limit}

For wavelengths much larger than the nearest-neighbor separation,
\begin{equation}
    kd_{\mathrm{nn}} \ll 1,
\end{equation}
one has
\begin{equation}
    \sin \left( \frac{kd_{\mathrm{nn}}}{2} \right) = \frac{kd_{\mathrm{nn}}}{2} + \mathcal{O}\left( k^3d_{\mathrm{nn}}^3 \right).
\end{equation}

The dispersion relation reduces to
\begin{equation}
    \omega^2 = c_{\mathrm{E}}^2\mathcal{C}^2k^2 + \mathcal{O}\left( k^4d_{\mathrm{nn}}^2 \right).
\end{equation}

Therefore,
\begin{equation}
    \lim_{k\rightarrow 0} \frac{\omega}{k} = \mathcal{C} c_{\mathrm{E}}.
\end{equation}

The continuum propagation speed is consequently
\begin{equation}
    \boxed{ c_P = \mathcal{C} c_{\mathrm{E}} }.
\end{equation}

Using the locally uniform endpoint factor gives
\begin{equation}
    \boxed{ c_P = c_{\mathrm{E}} \left( 1-\alpha_{\mathrm{G}} \right)^2 }.
\end{equation}

The constrained contact action therefore reproduces the propagation law introduced in Section 10.4.

The positive-frequency branch may be written as
\begin{equation}
    \omega(k,\alpha_{\mathrm{G}}) = \frac{2c_{\mathrm{E}}\mathcal{C}}{d_{\mathrm{nn}}} \left| \sin \left( \frac{kd_{\mathrm{nn}}}{2} \right) \right|.
\end{equation}

For wave numbers within the first monotonic branch, the corresponding group velocity is
\begin{equation}
    v_{\mathrm{g}} = \frac{\mathrm{d}\omega}{\mathrm{d}k} = c_{\mathrm{E}}\mathcal{C} \cos \left( \frac{kd_{\mathrm{nn}}}{2} \right).
\end{equation}

Hence,
\begin{equation}
    \lim_{k\rightarrow 0}v_{\mathrm{g}} = c_{\mathrm{E}}\mathcal{C} = c_P.
\end{equation}

The endpoint constraint therefore rescales both the long-wavelength phase velocity and the long-wavelength group velocity.

\subsection{Vector Form of the Transverse Contact Energy}

Let the unit vector along the FCC contact from node $i$ to node $j$ be
\begin{equation}
    \hat{\mathbf{n}}_{ij}.
\end{equation}

The projector onto the plane transverse to the contact is
\begin{equation}
    \mathbf{P}_{\perp}^{(ij)} = \mathbf{I} - \hat{\mathbf{n}}_{ij} \hat{\mathbf{n}}_{ij}^{\mathsf{T}}.
\end{equation}

For the relative nodal displacement
\begin{equation}
    \Delta\mathbf{u}_{ij} = \mathbf{u}_j - \mathbf{u}_i,
\end{equation}
the transverse contact energy per unit participating interface area is
\begin{equation}
    \boxed{ \mathcal{U}_{\perp}^{(ij)} = \frac{\tau_{\mathrm{D}}}{2d_{\mathrm{nn}}^2} \Delta\mathbf{u}_{ij}^{\mathsf{T}} \mathbf{P}_{\perp}^{(ij)} \Delta\mathbf{u}_{ij} }.
\end{equation}

The corresponding two-node block Hessian is
\begin{equation}
    \boxed{ \mathbf{H}_{\perp}^{(ij)} = \frac{\tau_{\mathrm{D}}}{d_{\mathrm{nn}}^2}
    \begin{pmatrix}
        \mathbf{P}_{\perp}^{(ij)} & -\mathbf{P}_{\perp}^{(ij)} \\[3pt]
        -\mathbf{P}_{\perp}^{(ij)} & \mathbf{P}_{\perp}^{(ij)}
    \end{pmatrix} }.
\end{equation}

The scalar dispersion derived above is the restriction of this vector contact law to one transverse polarization and one locally uniform propagation direction.

The full FCC dynamic matrix is obtained by summing the corresponding block contributions over the twelve nearest-neighbor directions of each node.

\subsection{Saturation Limit}

As the gravitational strain approaches the active-phase boundary,
\begin{equation}
    \alpha_{\mathrm{G}} \rightarrow 1,
\end{equation}
the endpoint-transfer factor vanishes:
\begin{equation}
    \mathcal{C} \rightarrow 0.
\end{equation}

The reduced surface inertia diverges,
\begin{equation}
    \mu_{\mathrm{eff}} = \frac{\mu_{\mathrm{D}}}{\mathcal{C}^2} \rightarrow \infty,
\end{equation}
while the transverse contact Hessian remains finite:
\begin{equation}
    \mathbf{H}_U(\alpha_{\mathrm{G}}) = \mathbf{H}_U(0).
\end{equation}

For every fixed lattice wave number,
\begin{equation}
    \omega(k,\alpha_{\mathrm{G}}) \rightarrow 0,
\end{equation}
and the long-wavelength propagation speed satisfies
\begin{equation}
    c_P \rightarrow 0.
\end{equation}

The active FCC mode is therefore arrested by divergent constrained inertia rather than by continuous loss of the baseline transverse stiffness.

This limiting behavior is defined only from below,
\begin{equation}
    \alpha_{\mathrm{G}} \rightarrow 1^{-}.
\end{equation}

At
\begin{equation}
    \alpha_{\mathrm{G}} = 1,
\end{equation}
the active FCC phase terminates and the local structure transitions to the phase-locked SC substrate. The divergent reduced inertia signals the loss of validity of the FCC dynamical coordinates; it is not interpreted as an infinite material mass within the SC phase.

\subsection{Constitutive Result and Scope}

The constrained transverse contact action yields
\begin{equation}
    \boxed{ \mathbf{H}_U(\alpha_{\mathrm{G}}) = \mathbf{H}_U(0) }
\end{equation}
and
\begin{equation}
    \boxed{ \mathbf{H}_T(\alpha_{\mathrm{G}}) = \frac{\mathbf{H}_T(0)}{ \left( 1-\alpha_{\mathrm{G}} \right)^4 } }.
\end{equation}

Its directional nearest-neighbor dispersion relation is
\begin{equation}
    \boxed{ \omega^2(k,\alpha_{\mathrm{G}}) = \frac{4c_{\mathrm{E}}^2}{d_{\mathrm{nn}}^2} \left( 1-\alpha_{\mathrm{G}} \right)^4 \sin^2 \left( \frac{kd_{\mathrm{nn}}}{2} \right) },
\end{equation}
with continuum limit
\begin{equation}
    \boxed{ \lim_{k\rightarrow 0} \frac{\omega}{k} = c_{\mathrm{E}} \left( 1-\alpha_{\mathrm{G}} \right)^2 }.
\end{equation}

The derivation establishes that the endpoint-transfer closure is compatible with a conservative transverse contact action and that its inertial and dispersive consequences are mutually consistent.

The result remains conditional on the local transfer constraint
\begin{equation}
    u_n = \mathcal{C} y_n, \qquad \mathcal{C} = \left[ 1-\alpha_{\mathrm{G}}^{(i)} \right] \left[ 1-\alpha_{\mathrm{G}}^{(j)} \right].
\end{equation}

A complete microscopic FCC constitutive theory must derive this transfer Jacobian from the normal, transverse, torsional, phase, condensate, and vertical-constraint sectors of the local contact energy. The present action supplies the required reduced limit that such a microscopic theory must recover.

\section{Reduced Constitutive Derivation of the Endpoint Transfer Jacobian}

\subsection{Constitutive Setting}

The preceding section introduced the local relation
\begin{equation}
    u_n = \mathcal{C}_{ij}y_n,
\end{equation}
where $y_n$ is the underlying condensate-supported reconfiguration amplitude, $u_n$ is the completed transverse displacement of the active FCC state, and $\mathcal{C}_{ij}$ is the transfer Jacobian associated with the two endpoints of a nearest-neighbor FCC contact.

The purpose of the present section is to derive $\mathcal{C}_{ij}$ from the local active-tension balance.

The analysis applies within the active FCC phase,
\begin{equation}
    0 \leq \alpha_{\mathrm{G}} < 1,
\end{equation}
and assumes that the gravitational background is quasistatic over the duration of a single nearest-neighbor transfer:
\begin{equation}
    \dot{\alpha}_{\mathrm{G}} \approx 0.
\end{equation}

The baseline surface shear tension $\tau_{\mathrm{D}}$ remains the tangent constitutive coefficient of the active Soft Floor. Gravitational loading does not continuously replace this coefficient with a softened modulus. Instead, it occupies a fraction of the finite active capacity available to support an additional reconfiguration.

This distinction separates:
\begin{equation}
    \tau_{\mathrm{D}},
\end{equation}
the baseline transverse constitutive tension of the active FCC phase, from
\begin{equation}
    \tau_{\mathrm{av}},
\end{equation}
the portion of that capacity remaining available after the local gravitational load has been applied.

\subsection{Residual Active Capacity at One Endpoint}

At endpoint $i$, the dimensionless gravitational strain is
\begin{equation}
    \alpha_{\mathrm{G}}^{(i)} = \frac{\tau_L^{(i)}}{\tau_{\mathrm{D}}},
\end{equation}
where $\tau_L^{(i)}$ is the local mass-induced loading of the Soft Floor.

The residual active tension capacity is therefore
\begin{equation}
    \tau_{\mathrm{av}}^{(i)} = \tau_{\mathrm{D}} - \tau_L^{(i)}.
\end{equation}

Substitution of the normalized loading gives
\begin{equation}
    \tau_{\mathrm{av}}^{(i)} = \tau_{\mathrm{D}} \left( 1-\alpha_{\mathrm{G}}^{(i)} \right).
\end{equation}

Hence,
\begin{equation}
    \boxed{ \frac{\tau_{\mathrm{av}}^{(i)}}{\tau_{\mathrm{D}}} = 1-\alpha_{\mathrm{G}}^{(i)} }.
\end{equation}

The same relation holds at endpoint $j$:
\begin{equation}
    \boxed{ \frac{\tau_{\mathrm{av}}^{(j)}}{\tau_{\mathrm{D}}} = 1-\alpha_{\mathrm{G}}^{(j)} }.
\end{equation}

These relations follow directly from the definition of $\alpha_{\mathrm{G}}$. No additional endpoint coefficient is introduced.

\subsection{Work-Conjugate Endpoint Coordinates}

Let $d_{\mathrm{nn}}$ be the nearest-neighbor separation of the active FCC lattice. For one selected transverse channel, define the dimensionless latent reconfiguration coordinate at endpoint $i$ by
\begin{equation}
    \eta_i = \frac{y_i}{d_{\mathrm{nn}}},
\end{equation}
and the dimensionless completed coordinate by
\begin{equation}
    \varepsilon_i = \frac{u_i}{d_{\mathrm{nn}}}.
\end{equation}

The two coordinates describe the same incremental endpoint event at different constitutive levels:
\begin{itemize}
    \item $\eta_i$ measures the underlying node--condensate reconfiguration;
    \item $\varepsilon_i$ measures the completed displacement transmitted into the active FCC state.
\end{itemize}

For an infinitesimal endpoint reconfiguration, the virtual work per unit participating interface area supplied by the remaining active capacity is
\begin{equation}
    \left( \frac{\delta W}{A} \right)_i = \tau_{\mathrm{av}}^{(i)} \delta\eta_i.
\end{equation}

The same completed event, expressed using the baseline transverse constitutive tension, is
\begin{equation}
    \left( \frac{\delta W}{A} \right)_i = \tau_{\mathrm{D}} \delta\varepsilon_i.
\end{equation}

Equality of the two work-conjugate representations requires
\begin{equation}
    \tau_{\mathrm{D}} \delta\varepsilon_i = \tau_{\mathrm{av}}^{(i)} \delta\eta_i.
\end{equation}

Using
\begin{equation}
    \tau_{\mathrm{av}}^{(i)} = \tau_{\mathrm{D}} \left( 1-\alpha_{\mathrm{G}}^{(i)} \right),
\end{equation}
gives
\begin{equation}
    \delta\varepsilon_i = \left( 1-\alpha_{\mathrm{G}}^{(i)} \right) \delta\eta_i.
\end{equation}

For a quasistatic local strain, $\alpha_{\mathrm{G}}^{(i)}$ is constant during the incremental transfer. Integration from the common undeformed reference state,
\begin{equation}
    \eta_i=0, \qquad \varepsilon_i=0,
\end{equation}
therefore yields
\begin{equation}
    \boxed{ \varepsilon_i = \left( 1-\alpha_{\mathrm{G}}^{(i)} \right) \eta_i }.
\end{equation}

Multiplying by $d_{\mathrm{nn}}$ gives the single-endpoint displacement map
\begin{equation}
    \boxed{ u_i = \left( 1-\alpha_{\mathrm{G}}^{(i)} \right) y_i }.
\end{equation}

The factor $(1-\alpha_{\mathrm{G}}^{(i)})$ is consequently the Jacobian between the latent reconfiguration coordinate and the completed active-FCC displacement at endpoint $i$.

This result does not represent dissipation. It is a conservative coordinate reduction produced by the finite active-capacity balance.

\subsection{Source Release}

Consider a nearest-neighbor transfer directed from node $i$ to node $j$. Let $y_{ij}$ denote the latent reconfiguration amplitude presented to the source endpoint.

The source endpoint releases only the fraction permitted by its residual active capacity. The released amplitude is therefore
\begin{equation}
    y_{ij}^{(\mathrm{rel})} = \left( 1-\alpha_{\mathrm{G}}^{(i)} \right) y_{ij}.
\end{equation}

The released state remains a coherent lattice reconfiguration. It becomes the incident state presented to the target endpoint.

\subsection{Target Acceptance and Closure}

The target endpoint must accept the incoming state and restore local FCC closure. Applying the single-endpoint map at node $j$ gives
\begin{equation}
    u_{ij} = \left( 1-\alpha_{\mathrm{G}}^{(j)} \right) y_{ij}^{(\mathrm{rel})}.
\end{equation}

Substitution of the source-release relation yields
\begin{equation}
    u_{ij} = \left( 1-\alpha_{\mathrm{G}}^{(j)} \right) \left( 1-\alpha_{\mathrm{G}}^{(i)} \right) y_{ij}.
\end{equation}

Therefore,
\begin{equation}
    \boxed{ u_{ij} = \mathcal{C}_{ij}y_{ij} },
\end{equation}
where
\begin{equation}
    \boxed{ \mathcal{C}_{ij} = \left[ 1-\alpha_{\mathrm{G}}^{(i)} \right] \left[ 1-\alpha_{\mathrm{G}}^{(j)} \right] }.
\end{equation}

The endpoint transfer Jacobian is thus the serial composition of the source-release map and the target-acceptance map.

\subsection{Origin of the Multiplicative Structure}

The two endpoint factors multiply because the corresponding constraints act sequentially on the same transferred state.

The local transfer sequence is
\begin{equation}
    \text{latent source state} \longrightarrow \text{source-released state} \longrightarrow \text{target-completed state}.
\end{equation}

The target endpoint does not act independently on the original latent amplitude. It acts on the amplitude that has already passed through the source constraint.

Consequently, the completed transfer is obtained by composition:
\begin{equation}
    y_{ij} \overset{ 1-\alpha_{\mathrm{G}}^{(i)} }{ \longrightarrow } y_{ij}^{(\mathrm{rel})} \overset{ 1-\alpha_{\mathrm{G}}^{(j)} }{ \longrightarrow } u_{ij}.
\end{equation}

The resulting transfer Jacobian is therefore a product rather than a sum or arithmetic average.

Because the endpoint maps are scalar within the selected transverse polarization,
\begin{equation}
    \mathcal{C}_{ij} = \mathcal{C}_{ji}.
\end{equation}

This reciprocity is required for a conservative nearest-neighbor contact in a static gravitational background.

\subsection{Limiting Behavior}

In the unstrained active vacuum,
\begin{equation}
    \alpha_{\mathrm{G}}^{(i)} = \alpha_{\mathrm{G}}^{(j)} = 0,
\end{equation}
and therefore
\begin{equation}
    \mathcal{C}_{ij}=1.
\end{equation}

The latent and completed coordinates coincide:
\begin{equation}
    u_{ij} = y_{ij}.
\end{equation}

If either endpoint approaches saturation,
\begin{equation}
    \alpha_{\mathrm{G}}^{(i)} \rightarrow 1 \qquad\text{or}\qquad \alpha_{\mathrm{G}}^{(j)} \rightarrow 1,
\end{equation}
then
\begin{equation}
    \mathcal{C}_{ij} \rightarrow 0.
\end{equation}

A completed active-FCC transfer cannot occur if either endpoint has exhausted its available reconfiguration capacity.

Throughout the active phase,
\begin{equation}
    0 < \mathcal{C}_{ij} \leq 1.
\end{equation}

At the exact saturation boundary, the FCC constitutive description terminates and the local structure enters the phase-locked SC state. The endpoint Jacobian is therefore defined only for contacts whose two endpoints remain in the active FCC phase.

\subsection{Local Continuum Reduction}

For a gravitational field varying negligibly across one nearest-neighbor separation,
\begin{equation}
    \alpha_{\mathrm{G}}^{(i)} \approx \alpha_{\mathrm{G}}^{(j)} \approx \alpha_{\mathrm{G}}.
\end{equation}

The endpoint Jacobian then reduces to
\begin{equation}
    \boxed{ \mathcal{C}_{ij} \approx \left( 1-\alpha_{\mathrm{G}} \right)^2 }.
\end{equation}

More generally, defining the endpoint mean and difference by
\begin{equation}
    \alpha_{\mathrm{G}} = \frac{ \alpha_{\mathrm{G}}^{(i)} + \alpha_{\mathrm{G}}^{(j)} }{2},
\end{equation}
and
\begin{equation}
    \delta\alpha_{\mathrm{G}} = \alpha_{\mathrm{G}}^{(j)} - \alpha_{\mathrm{G}}^{(i)},
\end{equation}
gives the exact identity
\begin{equation}
    \mathcal{C}_{ij} = \left( 1-\alpha_{\mathrm{G}} \right)^2 - \frac{ \left( \delta\alpha_{\mathrm{G}} \right)^2 }{4}.
\end{equation}

The first-order endpoint difference cancels. In the long-wavelength continuum regime,
\begin{equation}
    \left| \delta\alpha_{\mathrm{G}} \right| \ll 2 \left| 1-\alpha_{\mathrm{G}} \right|,
\end{equation}
and the quadratic correction is negligible. Hence,
\begin{equation}
    \boxed{ \mathcal{C} = \left( 1-\alpha_{\mathrm{G}} \right)^2 }.
\end{equation}

\subsection{Conservative Character of the Mapping}

The endpoint transformation is holonomic and time independent under the quasistatic-background condition. It therefore preserves the conservative structure of the local action.

For the underlying reconfiguration coordinate,
\begin{equation}
    T = \frac{1}{2} \mu_{\mathrm{D}}\dot{y}^2.
\end{equation}

Using
\begin{equation}
    u = \mathcal{C} y,
\end{equation}
with
\begin{equation}
    \dot{\mathcal{C}} = 0,
\end{equation}
gives
\begin{equation}
    \dot{y} = \frac{\dot{u}}{\mathcal{C}}.
\end{equation}

The kinetic energy becomes
\begin{equation}
    T = \frac{1}{2} \frac{\mu_{\mathrm{D}}}{\mathcal{C}^2} \dot{u}^2.
\end{equation}

No term linear in $\dot{u}$ is generated, and no dissipative coefficient is introduced. The reduced action remains invariant under time reversal.

The constrained part of the reconfiguration is retained in the increased kinetic loading required to produce a fixed completed displacement. It is not irreversibly converted into heat or lost from the lattice.

\subsection{Recovery of the Effective Inertia}

The endpoint Jacobian gives the effective surface inertia
\begin{equation}
    \boxed{ \mu_{\mathrm{eff}}^{(ij)} = \frac{\mu_{\mathrm{D}}}{\mathcal{C}_{ij}^2} }.
\end{equation}

In the local continuum limit,
\begin{equation}
    \mathcal{C} = \left( 1-\alpha_{\mathrm{G}} \right)^2,
\end{equation}
so
\begin{equation}
    \boxed{ \mu_{\mathrm{eff}} = \frac{\mu_{\mathrm{D}}}{ \left( 1-\alpha_{\mathrm{G}} \right)^4 } }.
\end{equation}

The inverse-fourth-power dependence therefore follows from two operations:
\begin{equation}
    \text{two serial endpoint factors} \longrightarrow \mathcal{C} = \left( 1-\alpha_{\mathrm{G}} \right)^2,
\end{equation}
followed by
\begin{equation}
    \text{quadratic kinetic reduction} \longrightarrow \mu_{\mathrm{eff}} = \frac{\mu_{\mathrm{D}}}{\mathcal{C}^2}.
\end{equation}

\subsection{Recovery of the Propagation Law}

The baseline transverse stiffness--inertia relation is
\begin{equation}
    c_{\mathrm{E}}^2 = \frac{\tau_{\mathrm{D}}}{\mu_{\mathrm{D}}}.
\end{equation}

The active-phase transverse potential Hessian remains unchanged, while the effective inertia becomes $\mu_{\mathrm{eff}}$. Therefore,
\begin{equation}
    c_P^2 = \frac{\tau_{\mathrm{D}}}{\mu_{\mathrm{eff}}}.
\end{equation}

Substitution gives
\begin{equation}
    c_P^2 = \frac{\tau_{\mathrm{D}}}{\mu_{\mathrm{D}}} \mathcal{C}^2.
\end{equation}

Using the baseline identity,
\begin{equation}
    c_P^2 = c_{\mathrm{E}}^2\mathcal{C}^2.
\end{equation}

Taking the positive propagation branch,
\begin{equation}
    c_P = c_{\mathrm{E}}\mathcal{C}.
\end{equation}

In the local continuum limit,
\begin{equation}
    \boxed{ c_P = c_{\mathrm{E}} \left( 1-\alpha_{\mathrm{G}} \right)^2 }.
\end{equation}

The gravitational propagation law is therefore the continuum consequence of the serial endpoint transfer Jacobian.

\subsection{Constitutive Result}

The endpoint-transfer Jacobian is
\begin{equation}
    \boxed{ \mathcal{C}_{ij} = \left[ 1-\alpha_{\mathrm{G}}^{(i)} \right] \left[ 1-\alpha_{\mathrm{G}}^{(j)} \right] }.
\end{equation}

It follows from:
\begin{equation}
    \alpha_{\mathrm{G}}^{(i)} = \frac{\tau_L^{(i)}}{\tau_{\mathrm{D}}},
\end{equation}
\begin{equation}
    \tau_{\mathrm{av}}^{(i)} = \tau_{\mathrm{D}} - \tau_L^{(i)},
\end{equation}
the work-conjugate equivalence between latent and completed endpoint coordinates, and the serial composition of source release with target acceptance.

The resulting constitutive chain is
\begin{equation}
    \boxed{ \alpha_{\mathrm{G}} \longrightarrow \tau_{\mathrm{av}} \longrightarrow \mathcal{C}_{ij} \longrightarrow \mu_{\mathrm{eff}} \longrightarrow c_P }.
\end{equation}

Explicitly,
\begin{equation}
    \boxed{ \mathcal{C} = \left( 1-\alpha_{\mathrm{G}} \right)^2 },
\end{equation}
\begin{equation}
    \boxed{ \mu_{\mathrm{eff}} = \frac{\mu_{\mathrm{D}}}{ \left( 1-\alpha_{\mathrm{G}} \right)^4 } },
\end{equation}
and
\begin{equation}
    \boxed{ c_P = c_{\mathrm{E}} \left( 1-\alpha_{\mathrm{G}} \right)^2 }.
\end{equation}

\subsection{Constitutive Scope}

The derivation establishes the endpoint Jacobian at the reduced microscopic contact level. It relies on the active-phase premises that:
\begin{equation}
    \alpha_{\mathrm{G}} = \frac{\tau_L}{\tau_{\mathrm{D}}}
\end{equation}
measures the fraction of the local transfer capacity occupied by gravitational loading, and that the remaining capacity performs conservative work through the unchanged baseline transverse constitutive tension.

A complete bond-level theory must recover this work-conjugate endpoint map from an explicit FCC contact free energy containing normal displacement, transverse shear, relative phase, condensate response, and vertical Hard-Floor constraint. The present result defines the reduced constitutive limit that such a contact law must reproduce.

\section{Objective FCC Contact Kinematics and Quadratic Bond Action}

\subsection{Scope}

The preceding sections derived the endpoint-transfer Jacobian and established its consequences for transverse propagation through a single nearest-neighbor direction. Extension to the complete FCC lattice requires a three-dimensional contact law that is invariant under rigid translation, rigid rotation, and uniform phase displacement.

A contact energy written solely in terms of the relative displacement of two node centers is insufficient for this purpose. Under a rigid-body rotation, neighboring node centers acquire a transverse relative displacement even though their physical contact surfaces undergo no elastic slip. The contact law must therefore include both nodal translations and nodal microrotations.

The present section develops the linearized objective kinematics and the corresponding parity-even quadratic action for one nearest-neighbor FCC contact. The construction applies to small incremental deformations about an active-phase equilibrium satisfying
\begin{equation}
    0 \leq \alpha_{\mathrm{G}} < 1.
\end{equation}

The contact coefficients introduced below are tangent constitutive coefficients. Their numerical reduction to the established Soft Floor constants is addressed through subsequent FCC homogenization.

\subsection{Reference Contact Geometry}

Consider two nearest-neighbor FCC nodes labeled $i$ and $j$. Their reference center-to-center vector is
\begin{equation}
    \mathbf{d}_{ij} = d_{\mathrm{nn}}\hat{\mathbf{n}}_{ij},
\end{equation}
where
\begin{equation}
    d_{\mathrm{nn}} = 2r_{\mathrm{D}}
\end{equation}
is the nearest-neighbor separation and
\begin{equation}
    \left| \hat{\mathbf{n}}_{ij} \right| = 1.
\end{equation}

The normal and transverse projectors associated with the contact are
\begin{equation}
    \mathbf{P}_{\mathrm{n}}^{(ij)} = \hat{\mathbf{n}}_{ij} \hat{\mathbf{n}}_{ij}^{\mathsf{T}},
\end{equation}
and
\begin{equation}
    \boxed{ \mathbf{P}_{\perp}^{(ij)} = \mathbf{I} - \hat{\mathbf{n}}_{ij} \hat{\mathbf{n}}_{ij}^{\mathsf{T}} }.
\end{equation}

Let
\begin{equation}
    \mathbf{u}_i, \qquad \mathbf{u}_j
\end{equation}
denote the infinitesimal translations of the node centers.

Let
\begin{equation}
    \boldsymbol{\vartheta}_i, \qquad \boldsymbol{\vartheta}_j
\end{equation}
denote their infinitesimal microrotations. These variables describe rotation of the local contact frames and are distinct from the internal torsional phase of the Soft Floor.

The internal phases of the two nodes are denoted by
\begin{equation}
    \varphi_i, \qquad \varphi_j.
\end{equation}

The microrotation $\boldsymbol{\vartheta}$ is a three-dimensional mechanical rotation of the node surface. The scalar $\varphi$ is the internal phase variable governing topological and torsional closure. The two degrees of freedom are not identified at the present constitutive order.

\subsection{Objective Relative Contact Displacement}

The reference contact point lies at
\begin{equation}
    +r_{\mathrm{D}}\hat{\mathbf{n}}_{ij}
\end{equation}
relative to the center of node $i$, and at
\begin{equation}
    -r_{\mathrm{D}}\hat{\mathbf{n}}_{ij}
\end{equation}
relative to the center of node $j$.

The incremental displacement of the contact point carried by node $i$ is therefore
\begin{equation}
    \mathbf{u}_i^{\mathrm{c}} = \mathbf{u}_i + r_{\mathrm{D}} \boldsymbol{\vartheta}_i \times \hat{\mathbf{n}}_{ij}.
\end{equation}

The corresponding displacement of the contact point carried by node $j$ is
\begin{equation}
    \mathbf{u}_j^{\mathrm{c}} = \mathbf{u}_j - r_{\mathrm{D}} \boldsymbol{\vartheta}_j \times \hat{\mathbf{n}}_{ij}.
\end{equation}

The relative displacement of the two material contact points is
\begin{equation}
\begin{aligned}
    \Delta\mathbf{u}_{ij}^{\mathrm{c}} &= \mathbf{u}_j^{\mathrm{c}} - \mathbf{u}_i^{\mathrm{c}} \\[3pt]
    &= \mathbf{u}_j - \mathbf{u}_i - r_{\mathrm{D}} \left( \boldsymbol{\vartheta}_i + \boldsymbol{\vartheta}_j \right) \times \hat{\mathbf{n}}_{ij}.
\end{aligned}
\end{equation}

Thus,
\begin{equation}
    \boxed{ \Delta\mathbf{u}_{ij}^{\mathrm{c}} = \mathbf{u}_j - \mathbf{u}_i - r_{\mathrm{D}} \left( \boldsymbol{\vartheta}_i + \boldsymbol{\vartheta}_j \right) \times \hat{\mathbf{n}}_{ij} }.
\end{equation}

This contact-point displacement, rather than the relative center displacement alone, is the objective kinematic quantity entering the local shear energy.

\subsection{Normal Contact Deformation}

The normal contact deformation is the projection of the relative contact displacement onto the bond direction:
\begin{equation}
    \delta_{\mathrm{n}}^{(ij)} = \hat{\mathbf{n}}_{ij} \cdot \Delta\mathbf{u}_{ij}^{\mathrm{c}}.
\end{equation}

The rotational contribution is transverse to $\hat{\mathbf{n}}_{ij}$, and therefore
\begin{equation}
    \boxed{ \delta_{\mathrm{n}}^{(ij)} = \hat{\mathbf{n}}_{ij} \cdot \left( \mathbf{u}_j - \mathbf{u}_i \right) }.
\end{equation}

Positive and negative values represent normal separation and compression relative to the selected reference convention.

\subsection{Tangential Contact Slip}

The objective tangential slip is
\begin{equation}
    \boxed{ \boldsymbol{\delta}_{\mathrm{t}}^{(ij)} = \mathbf{P}_{\perp}^{(ij)} \Delta\mathbf{u}_{ij}^{\mathrm{c}} }.
\end{equation}

Explicitly,
\begin{equation}
    \boxed{ \boldsymbol{\delta}_{\mathrm{t}}^{(ij)} = \mathbf{P}_{\perp}^{(ij)} \left[ \mathbf{u}_j - \mathbf{u}_i - r_{\mathrm{D}} \left( \boldsymbol{\vartheta}_i + \boldsymbol{\vartheta}_j \right) \times \hat{\mathbf{n}}_{ij} \right] }.
\end{equation}

This quantity measures the relative displacement of the two physical contact surfaces in the plane tangent to the bond.

It differs fundamentally from
\begin{equation}
    \mathbf{P}_{\perp}^{(ij)} \left( \mathbf{u}_j - \mathbf{u}_i \right),
\end{equation}
which does not remove the contact-point motion produced by a rigid rotation of the pair.

\subsection{Relative Microrotation}

Define the relative microrotation by
\begin{equation}
    \Delta\boldsymbol{\vartheta}_{ij} = \boldsymbol{\vartheta}_j - \boldsymbol{\vartheta}_i.
\end{equation}

Its transverse component describes differential rolling or bending of the two contact frames:
\begin{equation}
    \boxed{ \boldsymbol{\delta}_{\mathrm{b}}^{(ij)} = \mathbf{P}_{\perp}^{(ij)} \Delta\boldsymbol{\vartheta}_{ij} }.
\end{equation}

Its component along the bond direction describes relative twist:
\begin{equation}
    \boxed{ \delta_{\mathrm{tw}}^{(ij)} = \hat{\mathbf{n}}_{ij} \cdot \Delta\boldsymbol{\vartheta}_{ij} }.
\end{equation}

Both quantities vanish when the two nodes undergo the same rigid microrotation.

\subsection{Internal Phase Difference}

The contact phase difference is
\begin{equation}
    \boxed{ \Delta\varphi_{ij} = \varphi_j - \varphi_i }.
\end{equation}

Only the phase difference can enter the contact energy. Consequently, the bond action is invariant under a uniform phase displacement,
\begin{equation}
    \varphi_i \rightarrow \varphi_i + \varphi_0, \qquad \varphi_j \rightarrow \varphi_j + \varphi_0.
\end{equation}

This invariance is the local contact expression of the global phase freedom of the Soft Floor.

\subsection{Linearized Frame Indifference}

Consider an arbitrary infinitesimal rigid-body motion consisting of a translation $\mathbf{a}$ and a rotation $\boldsymbol{\Omega}$. The nodal displacements are
\begin{equation}
    \mathbf{u}_i = \mathbf{a} + \boldsymbol{\Omega} \times \mathbf{X}_i,
\end{equation}
and
\begin{equation}
    \mathbf{u}_j = \mathbf{a} + \boldsymbol{\Omega} \times \mathbf{X}_j,
\end{equation}
where $\mathbf{X}_i$ and $\mathbf{X}_j$ are the reference node positions.

Since
\begin{equation}
    \mathbf{X}_j - \mathbf{X}_i = d_{\mathrm{nn}} \hat{\mathbf{n}}_{ij},
\end{equation}
the relative center displacement is
\begin{equation}
    \mathbf{u}_j - \mathbf{u}_i = d_{\mathrm{nn}} \boldsymbol{\Omega} \times \hat{\mathbf{n}}_{ij}.
\end{equation}

For a rigid rotation of the node frames,
\begin{equation}
    \boldsymbol{\vartheta}_i = \boldsymbol{\vartheta}_j = \boldsymbol{\Omega}.
\end{equation}

The relative contact-point displacement becomes
\begin{equation}
\begin{aligned}
    \Delta\mathbf{u}_{ij}^{\mathrm{c}} &= d_{\mathrm{nn}} \boldsymbol{\Omega} \times \hat{\mathbf{n}}_{ij} - 2r_{\mathrm{D}} \boldsymbol{\Omega} \times \hat{\mathbf{n}}_{ij}.
\end{aligned}
\end{equation}

Using
\begin{equation}
    d_{\mathrm{nn}} = 2r_{\mathrm{D}},
\end{equation}
one obtains
\begin{equation}
    \boxed{ \Delta\mathbf{u}_{ij}^{\mathrm{c}} = \mathbf{0} }.
\end{equation}

It follows that
\begin{equation}
    \delta_{\mathrm{n}}^{(ij)} = 0,
\end{equation}
and
\begin{equation}
    \boldsymbol{\delta}_{\mathrm{t}}^{(ij)} = \mathbf{0}.
\end{equation}

Because the two nodes undergo the same microrotation,
\begin{equation}
    \boldsymbol{\delta}_{\mathrm{b}}^{(ij)} = \mathbf{0}, \qquad \delta_{\mathrm{tw}}^{(ij)} = 0.
\end{equation}

A uniform phase displacement also gives
\begin{equation}
    \Delta\varphi_{ij} = 0.
\end{equation}

The complete contact energy therefore vanishes under rigid translation, rigid rotation, and uniform phase displacement. The bond law is consequently frame-indifferent to the order retained in the quadratic theory.

\subsection{Quadratic Contact Energy}

The local contact energy is written per unit participating interface area. Introduce the tangent coefficients
\begin{equation}
    \kappa_{\mathrm{n}}, \qquad \kappa_{\mathrm{t}}, \qquad \kappa_{\mathrm{b}}, \qquad \kappa_{\mathrm{tw}}, \qquad \kappa_{\varphi},
\end{equation}
together with the tangential--bending coupling coefficient
\begin{equation}
    \kappa_{\mathrm{tb}}.
\end{equation}

All six coefficients are defined with dimensions
\begin{equation}
    [\kappa] = \mathrm{N/m^3},
\end{equation}
so that rotational and phase coordinates enter through the corresponding arc-length scales $r_{\mathrm{D}}\boldsymbol{\vartheta}$ and $r_{\mathrm{D}}\varphi$.

The most general parity-even quadratic energy retained at the present contact order is
\begin{equation}
    \boxed{
    \begin{aligned}
        \mathcal{U}_{ij}^{(2)} ={}& \frac{1}{2} \kappa_{\mathrm{n}} \left( \delta_{\mathrm{n}}^{(ij)} \right)^2 + \frac{1}{2} \kappa_{\mathrm{t}} \left| \boldsymbol{\delta}_{\mathrm{t}}^{(ij)} \right|^2 \\[3pt]
        &+ \kappa_{\mathrm{tb}}r_{\mathrm{D}} \boldsymbol{\delta}_{\mathrm{t}}^{(ij)} \cdot \boldsymbol{\delta}_{\mathrm{b}}^{(ij)} + \frac{1}{2} \kappa_{\mathrm{b}}r_{\mathrm{D}}^2 \left| \boldsymbol{\delta}_{\mathrm{b}}^{(ij)} \right|^2 \\[3pt]
        &+ \frac{1}{2} \kappa_{\mathrm{tw}}r_{\mathrm{D}}^2 \left( \delta_{\mathrm{tw}}^{(ij)} \right)^2 + \frac{1}{2} \kappa_{\varphi}r_{\mathrm{D}}^2 \left( \Delta\varphi_{ij} \right)^2.
    \end{aligned}
    }
\end{equation}

The tangential--bending term is permitted because both
\begin{equation}
    \boldsymbol{\delta}_{\mathrm{t}}^{(ij)}
\end{equation}
and
\begin{equation}
    r_{\mathrm{D}} \boldsymbol{\delta}_{\mathrm{b}}^{(ij)}
\end{equation}
transform as transverse vectors under the local axial symmetry of the contact.

The uncoupled contact approximation is recovered by setting
\begin{equation}
    \kappa_{\mathrm{tb}} = 0.
\end{equation}

The phase contribution may be extended beyond the harmonic regime by the periodic completion
\begin{equation}
    \boxed{ \mathcal{U}_{\varphi}^{(ij)} = \kappa_{\varphi}r_{\mathrm{D}}^2 \left[ 1 - \cos \left( \Delta\varphi_{ij} \right) \right] }.
\end{equation}

For
\begin{equation}
    \left| \Delta\varphi_{ij} \right| \ll 1,
\end{equation}
this reduces to
\begin{equation}
    \mathcal{U}_{\varphi}^{(ij)} = \frac{1}{2} \kappa_{\varphi}r_{\mathrm{D}}^2 \left( \Delta\varphi_{ij} \right)^2 + \mathcal{O}\left( \Delta\varphi_{ij}^4 \right).
\end{equation}

Thus, the quadratic phase term is the tangent Hessian of an exactly $2\pi$-periodic contact potential.

\subsection{Positive-Stability Conditions}

The normal, twist, and phase sectors require
\begin{equation}
    \kappa_{\mathrm{n}} > 0, \qquad \kappa_{\mathrm{tw}} > 0, \qquad \kappa_{\varphi} > 0.
\end{equation}

The coupled tangential--bending block is positive definite when
\begin{equation}
    \kappa_{\mathrm{t}} > 0, \qquad \kappa_{\mathrm{b}} > 0,
\end{equation}
and
\begin{equation}
    \boxed{ \kappa_{\mathrm{t}}\kappa_{\mathrm{b}} - \kappa_{\mathrm{tb}}^2 > 0 }.
\end{equation}

These inequalities ensure that the quadratic contact energy is non-negative and that no local deformation mode possesses negative incremental stiffness.

\subsection{Compact Contact-Strain Representation}

Define the local contact-strain vector
\begin{equation}
    \boldsymbol{\xi}_{ij} =
    \begin{pmatrix}
        \delta_{\mathrm{n}}^{(ij)} \\
        \boldsymbol{\delta}_{\mathrm{t}}^{(ij)} \\
        r_{\mathrm{D}}\boldsymbol{\delta}_{\mathrm{b}}^{(ij)} \\
        r_{\mathrm{D}}\delta_{\mathrm{tw}}^{(ij)} \\
        r_{\mathrm{D}}\Delta\varphi_{ij}
    \end{pmatrix}.
\end{equation}

The contact energy can then be written as
\begin{equation}
    \boxed{ \mathcal{U}_{ij}^{(2)} = \frac{1}{2} \boldsymbol{\xi}_{ij}^{\mathsf{T}} \mathbf{H}_{\mathrm{c}} \boldsymbol{\xi}_{ij} },
\end{equation}
where the contact tangent matrix has the block structure
\begin{equation}
    \boxed{ \mathbf{H}_{\mathrm{c}} =
    \begin{pmatrix}
        \kappa_{\mathrm{n}} & 0 & 0 & 0 & 0 \\
        0 & \kappa_{\mathrm{t}}\mathbf{I}_2 & \kappa_{\mathrm{tb}}\mathbf{I}_2 & 0 & 0 \\
        0 & \kappa_{\mathrm{tb}}\mathbf{I}_2 & \kappa_{\mathrm{b}}\mathbf{I}_2 & 0 & 0 \\
        0 & 0 & 0 & \kappa_{\mathrm{tw}} & 0 \\
        0 & 0 & 0 & 0 & \kappa_{\varphi}
    \end{pmatrix} }.
\end{equation}

The two-dimensional identity $\mathbf{I}_2$ acts within the contact tangent plane.

Let the generalized nodal coordinate vector be
\begin{equation}
    \mathbf{q}_{ij} =
    \begin{pmatrix}
        \mathbf{u}_i \\
        \boldsymbol{\vartheta}_i \\
        \varphi_i \\
        \mathbf{u}_j \\
        \boldsymbol{\vartheta}_j \\
        \varphi_j
    \end{pmatrix}.
\end{equation}

Because every component of $\boldsymbol{\xi}_{ij}$ is linear in the generalized coordinates,
\begin{equation}
    \boldsymbol{\xi}_{ij} = \mathbf{B}_{ij}\mathbf{q}_{ij},
\end{equation}
where $\mathbf{B}_{ij}$ is the objective bond-kinematic operator.

The quadratic contact energy becomes
\begin{equation}
    \mathcal{U}_{ij}^{(2)} = \frac{1}{2} \mathbf{q}_{ij}^{\mathsf{T}} \mathbf{K}_{ij} \mathbf{q}_{ij},
\end{equation}
with bond stiffness matrix
\begin{equation}
    \boxed{ \mathbf{K}_{ij} = \mathbf{B}_{ij}^{\mathsf{T}} \mathbf{H}_{\mathrm{c}} \mathbf{B}_{ij} }.
\end{equation}

This representation supplies the local stiffness block required for the subsequent twelve-bond FCC assembly.

\subsection{Reciprocal Contact Symmetry}

Interchanging the endpoint labels gives
\begin{equation}
    i \leftrightarrow j, \qquad \hat{\mathbf{n}}_{ij} \rightarrow -\hat{\mathbf{n}}_{ij}.
\end{equation}

Under this transformation,
\begin{equation}
    \boldsymbol{\delta}_{\mathrm{t}}^{(ij)} \rightarrow -\boldsymbol{\delta}_{\mathrm{t}}^{(ij)},
\end{equation}
and
\begin{equation}
    \boldsymbol{\delta}_{\mathrm{b}}^{(ij)} \rightarrow -\boldsymbol{\delta}_{\mathrm{b}}^{(ij)}.
\end{equation}

Their scalar product is unchanged. The normal and twist measures are also unchanged, while
\begin{equation}
    \Delta\varphi_{ij} \rightarrow -\Delta\varphi_{ij}.
\end{equation}

Because the phase energy is even in $\Delta\varphi_{ij}$, the complete bond energy satisfies
\begin{equation}
    \boxed{ \mathcal{U}_{ij} = \mathcal{U}_{ji} }.
\end{equation}

The bond is therefore reciprocal in a static, unpolarized Soft Floor background.

Parity-odd contact invariants are excluded from the present baseline. Such terms may become active in a locally chiral topological-defect sector, but they cancel in the unpolarized vacuum containing equal left- and right-handed closure senses.

\subsection{Constrained Kinetic Action}

Let
\begin{equation}
    \mathbf{y}_{ij}
\end{equation}
denote the latent generalized contact coordinate supported by the node--condensate system, and let
\begin{equation}
    \mathbf{q}_{ij}
\end{equation}
denote the corresponding completed FCC coordinate.

The endpoint-transfer relation is represented by
\begin{equation}
    \boxed{ \mathbf{q}_{ij} = \mathbf{T}_{ij}\mathbf{y}_{ij} }.
\end{equation}

The transfer operator acts as the previously derived endpoint Jacobian on the transverse state-transfer sector:
\begin{equation}
    \mathcal{C}_{ij} = \left[ 1-\alpha_{\mathrm{G}}^{(i)} \right] \left[ 1-\alpha_{\mathrm{G}}^{(j)} \right].
\end{equation}

Let
\begin{equation}
    \mathbb{P}_{\mathrm{tr}}^{(ij)}
\end{equation}
denote the generalized projector onto that transverse transfer sector. The local transfer operator is
\begin{equation}
    \boxed{ \mathbf{T}_{ij} = \mathbf{I} + \left( \mathcal{C}_{ij}-1 \right) \mathbb{P}_{\mathrm{tr}}^{(ij)} }.
\end{equation}

The complementary normal, microrotational, and internal-phase sectors remain unscaled at this constitutive order unless an additional loading dependence is derived for them.

Let the unstrained generalized kinetic energy per unit participating interface area be
\begin{equation}
    \mathcal{T}_{ij}^{(0)} = \frac{1}{2} \dot{\mathbf{y}}_{ij}^{\mathsf{T}} \mathbf{M}_{ij}^{(0)} \dot{\mathbf{y}}_{ij},
\end{equation}
where
\begin{equation}
    \mathbf{M}_{ij}^{(0)}
\end{equation}
is the symmetric positive inertial matrix of the unstrained node--condensate contact.

For a quasistatic gravitational background,
\begin{equation}
    \dot{\mathbf{T}}_{ij} = 0.
\end{equation}

Therefore,
\begin{equation}
    \dot{\mathbf{y}}_{ij} = \mathbf{T}_{ij}^{-1} \dot{\mathbf{q}}_{ij}.
\end{equation}

The reduced kinetic energy is
\begin{equation}
    \mathcal{T}_{ij} = \frac{1}{2} \dot{\mathbf{q}}_{ij}^{\mathsf{T}} \mathbf{M}_{ij}^{(\mathrm{eff})} \dot{\mathbf{q}}_{ij},
\end{equation}
with
\begin{equation}
    \boxed{ \mathbf{M}_{ij}^{(\mathrm{eff})} = \mathbf{T}_{ij}^{-\mathsf{T}} \mathbf{M}_{ij}^{(0)} \mathbf{T}_{ij}^{-1} }.
\end{equation}

On the transverse transfer sector,
\begin{equation}
    \boxed{ \mathbb{P}_{\mathrm{tr}}^{(ij)} \mathbf{M}_{ij}^{(\mathrm{eff})} \mathbb{P}_{\mathrm{tr}}^{(ij)} = \frac{1}{\mathcal{C}_{ij}^2} \mathbb{P}_{\mathrm{tr}}^{(ij)} \mathbf{M}_{ij}^{(0)} \mathbb{P}_{\mathrm{tr}}^{(ij)} }.
\end{equation}

For a locally uniform gravitational strain,
\begin{equation}
    \mathcal{C}_{ij} \approx \left( 1-\alpha_{\mathrm{G}} \right)^2,
\end{equation}
and the transverse inertial block becomes
\begin{equation}
    \boxed{ \mathbf{M}_{\mathrm{tr}}^{(\mathrm{eff})} = \frac{ \mathbf{M}_{\mathrm{tr}}^{(0)} }{ \left( 1-\alpha_{\mathrm{G}} \right)^4 } }.
\end{equation}

The potential Hessian is not multiplied by this factor:
\begin{equation}
    \boxed{ \mathbf{K}_{ij}(\alpha_{\mathrm{G}}) = \mathbf{K}_{ij}(0) }, \qquad 0 \leq \alpha_{\mathrm{G}} < 1.
\end{equation}

The gravitational strain therefore enters the objective contact action through the constrained kinetic map rather than through continuous softening of the active-phase contact Hessian.

\subsection{Quadratic Bond Action}

The reduced quadratic action for one nearest-neighbor contact is
\begin{equation}
    \boxed{ S_{ij}^{(2)} = \int \left[ \frac{1}{2} \dot{\mathbf{q}}_{ij}^{\mathsf{T}} \mathbf{M}_{ij}^{(\mathrm{eff})} \dot{\mathbf{q}}_{ij} - \frac{1}{2} \mathbf{q}_{ij}^{\mathsf{T}} \mathbf{K}_{ij} \mathbf{q}_{ij} \right] \mathrm{d}t }.
\end{equation}

The associated linearized contact equation is
\begin{equation}
    \boxed{ \mathbf{M}_{ij}^{(\mathrm{eff})} \ddot{\mathbf{q}}_{ij} + \mathbf{K}_{ij} \mathbf{q}_{ij} = \mathbf{0} }.
\end{equation}

This is the local generalized eigenvalue structure from which the FCC dynamic matrix is constructed.

\subsection{Null-Mode Conditions}

The contact action satisfies the three required symmetry null modes.

For uniform translation,
\begin{equation}
    \mathbf{u}_i = \mathbf{u}_j = \mathbf{a}, \qquad \boldsymbol{\vartheta}_i = \boldsymbol{\vartheta}_j = \mathbf{0}, \qquad \varphi_i = \varphi_j,
\end{equation}
one has
\begin{equation}
    \boldsymbol{\xi}_{ij} = \mathbf{0}.
\end{equation}

For an infinitesimal rigid rotation,
\begin{equation}
    \mathbf{u}_j - \mathbf{u}_i = d_{\mathrm{nn}} \boldsymbol{\Omega} \times \hat{\mathbf{n}}_{ij},
\end{equation}
and
\begin{equation}
    \boldsymbol{\vartheta}_i = \boldsymbol{\vartheta}_j = \boldsymbol{\Omega},
\end{equation}
one again obtains
\begin{equation}
    \boldsymbol{\xi}_{ij} = \mathbf{0}.
\end{equation}

For a uniform internal phase displacement,
\begin{equation}
    \varphi_i \rightarrow \varphi_i + \varphi_0, \qquad \varphi_j \rightarrow \varphi_j + \varphi_0,
\end{equation}
the phase difference remains unchanged.

Therefore,
\begin{equation}
    \boxed{ \mathbf{K}_{ij} \mathbf{q}_{\mathrm{rigid}} = \mathbf{0} }
\end{equation}
for all rigid translations, rigid rotations, and uniform phase shifts.

These null-space properties must be preserved when the twelve contact blocks are assembled into the complete FCC unit-cell Hessian.

\subsection{Recovery of the Directional Transverse Contact Law}

For a mode with no normal deformation, no relative microrotation, and no internal phase difference,
\begin{equation}
    \delta_{\mathrm{n}}^{(ij)} = 0,
\end{equation}
\begin{equation}
    \boldsymbol{\delta}_{\mathrm{b}}^{(ij)} = \mathbf{0},
\end{equation}
\begin{equation}
    \delta_{\mathrm{tw}}^{(ij)} = 0,
\end{equation}
and
\begin{equation}
    \Delta\varphi_{ij} = 0.
\end{equation}

The contact energy reduces to
\begin{equation}
    \mathcal{U}_{ij}^{(2)} = \frac{1}{2} \kappa_{\mathrm{t}} \left| \boldsymbol{\delta}_{\mathrm{t}}^{(ij)} \right|^2.
\end{equation}

For the restricted directional mode used in the preceding transverse-chain derivation, the normalization is
\begin{equation}
    \boxed{ \kappa_{\mathrm{t}} = \frac{\tau_{\mathrm{D}}}{d_{\mathrm{nn}}^2} }.
\end{equation}

Hence,
\begin{equation}
    \mathcal{U}_{ij}^{(2)} = \frac{\tau_{\mathrm{D}}}{2d_{\mathrm{nn}}^2} \left| \boldsymbol{\delta}_{\mathrm{t}}^{(ij)} \right|^2,
\end{equation}
which recovers the previously established scalar transverse contact energy as a special polarization of the objective three-dimensional law.

\subsection{Constitutive Result}

The objective nearest-neighbor contact is characterized by the deformation set
\begin{equation}
    \boxed{ \left\{ \delta_{\mathrm{n}}^{(ij)}, \; \boldsymbol{\delta}_{\mathrm{t}}^{(ij)}, \; \boldsymbol{\delta}_{\mathrm{b}}^{(ij)}, \; \delta_{\mathrm{tw}}^{(ij)}, \; \Delta\varphi_{ij} \right\} }.
\end{equation}

The corresponding stiffness matrix is
\begin{equation}
    \boxed{ \mathbf{K}_{ij} = \mathbf{B}_{ij}^{\mathsf{T}} \mathbf{H}_{\mathrm{c}} \mathbf{B}_{ij} },
\end{equation}
and the gravitationally constrained inertial matrix is
\begin{equation}
    \boxed{ \mathbf{M}_{ij}^{(\mathrm{eff})} = \mathbf{T}_{ij}^{-\mathsf{T}} \mathbf{M}_{ij}^{(0)} \mathbf{T}_{ij}^{-1} }.
\end{equation}

Within the transverse transfer sector,
\begin{equation}
    \boxed{ \mathbf{M}_{\mathrm{tr}}^{(\mathrm{eff})} = \frac{ \mathbf{M}_{\mathrm{tr}}^{(0)} }{ \mathcal{C}_{ij}^2 } },
\end{equation}
where
\begin{equation}
    \boxed{ \mathcal{C}_{ij} = \left[ 1-\alpha_{\mathrm{G}}^{(i)} \right] \left[ 1-\alpha_{\mathrm{G}}^{(j)} \right] }.
\end{equation}

The local contact theory therefore combines:
\begin{equation}
    \text{objective FCC surface kinematics},
\end{equation}
\begin{equation}
    \text{finite active-phase contact stiffness},
\end{equation}
and
\begin{equation}
    \text{gravitationally amplified constrained inertia}.
\end{equation}

The action supplies the individual bond block required for assembly over the twelve nearest-neighbor FCC directions. It does not yet impose continuum isotropy or assign final numerical values to the normal, rolling, twisting, and phase stiffnesses. Those quantities are determined by the subsequent lattice homogenization and matching to the established Soft Floor constitutive constants.

\section{Twelve-Bond FCC Dynamic Matrix and Long-Wavelength Homogenization}

\subsection{Scope}

The preceding section established an objective quadratic action for one nearest-neighbor contact of the active FCC Soft Floor. The present section assembles that contact law over the complete nearest-neighbor environment and derives the leading long-wavelength elastic tensor of the active vacuum.

The analysis is first performed in the unstrained active phase,
\begin{equation}
    \alpha_{\mathrm{G}} = 0,
\end{equation}
so that the lattice geometry and contact constitutive response are determined independently of gravitational loading. The endpoint-transfer Jacobian is introduced only after the unstrained continuum tensor has been obtained.

The purposes of the homogenization are to determine:
\begin{equation}
    C_{11}, \qquad C_{12}, \qquad C_{44},
\end{equation}
to identify the condition under which the FCC contact network is macroscopically isotropic, and to verify that the three-dimensional transverse acoustic branches retain the propagation law
\begin{equation}
    c_P = c_{\mathrm{E}} \left( 1-\alpha_{\mathrm{G}} \right)^2.
\end{equation}

The normal, tangential, microrotational, and internal-phase contact sectors are retained in the discrete action. At leading affine order, only those sectors that contribute to the classical elastic tensor remain active; the other sectors govern finite-wave-number and higher-gradient lattice dynamics.

\subsection{FCC Nearest-Neighbor Geometry}

The active Soft Floor is an FCC Bravais lattice with one generalized node per primitive cell. Each node has twelve directed nearest-neighbor connections. Because opposite directions describe the same undirected contact, the contact energy of one primitive cell is assembled from six independently counted bonds.

The nearest-neighbor separation is
\begin{equation}
    d_{\mathrm{nn}} = 2r_{\mathrm{D}}.
\end{equation}

A convenient set of six undirected nearest-neighbor vectors is
\begin{equation}
    \mathbf{d}_1 = \frac{d_{\mathrm{nn}}}{\sqrt{2}} (1,1,0),
\end{equation}
\begin{equation}
    \mathbf{d}_2 = \frac{d_{\mathrm{nn}}}{\sqrt{2}} (1,-1,0),
\end{equation}
\begin{equation}
    \mathbf{d}_3 = \frac{d_{\mathrm{nn}}}{\sqrt{2}} (1,0,1),
\end{equation}
\begin{equation}
    \mathbf{d}_4 = \frac{d_{\mathrm{nn}}}{\sqrt{2}} (1,0,-1),
\end{equation}
\begin{equation}
    \mathbf{d}_5 = \frac{d_{\mathrm{nn}}}{\sqrt{2}} (0,1,1),
\end{equation}
and
\begin{equation}
    \mathbf{d}_6 = \frac{d_{\mathrm{nn}}}{\sqrt{2}} (0,1,-1).
\end{equation}

The associated unit vectors are
\begin{equation}
    \hat{\mathbf{n}}_a = \frac{\mathbf{d}_a}{d_{\mathrm{nn}}}, \qquad a=1,\ldots,6.
\end{equation}

The negatives of these vectors supply the remaining six directed neighbor orientations.

The conventional cubic cell has edge length
\begin{equation}
    a_{\mathrm{FCC}} = 2\sqrt{2}r_{\mathrm{D}},
\end{equation}
and contains four lattice nodes. The primitive-cell volume is therefore
\begin{equation}
    \boxed{ V_{\mathrm{p}} = \frac{a_{\mathrm{FCC}}^3}{4} = 4\sqrt{2}r_{\mathrm{D}}^3 }.
\end{equation}

The six independently counted bonds per primitive cell are the discrete contact set used in the homogenized energy.

\subsection{Generalized Primitive-Cell Coordinate}

The generalized coordinate of the primitive cell at lattice position $\mathbf{R}$ is
\begin{equation}
    \mathbf{q}_{\mathbf{R}} =
    \begin{pmatrix}
        \mathbf{u}_{\mathbf{R}} \\
        \boldsymbol{\vartheta}_{\mathbf{R}} \\
        \varphi_{\mathbf{R}}
    \end{pmatrix},
\end{equation}
where:
\begin{itemize}
    \item $\mathbf{u}_{\mathbf{R}}$ is the nodal translation;
    \item $\boldsymbol{\vartheta}_{\mathbf{R}}$ is the nodal microrotation;
    \item $\varphi_{\mathbf{R}}$ is the internal torsional phase.
\end{itemize}

For a Bloch mode,
\begin{equation}
    \mathbf{q}_{\mathbf{R}}(t) = \mathbf{q}_0 \exp \left[ i \left( \mathbf{k}\cdot\mathbf{R}-\omega t \right) \right].
\end{equation}

The neighboring cell connected through $\mathbf{d}_a$ therefore satisfies
\begin{equation}
    \mathbf{q}_{\mathbf{R}+\mathbf{d}_a} = e^{i\mathbf{k}\cdot\mathbf{d}_a} \mathbf{q}_{\mathbf{R}}.
\end{equation}

For contact $a$, the objective contact-strain vector developed in the preceding section can be expressed as
\begin{equation}
    \boldsymbol{\xi}_a = \mathbf{B}_a(\mathbf{k})\mathbf{q}_0,
\end{equation}
where $\mathbf{B}_a(\mathbf{k})$ is the corresponding Bloch contact-kinematic operator.

Let $\mathbf{H}_{\mathrm{c}}$ denote the local contact tangent matrix. The stiffness contribution of contact $a$ is
\begin{equation}
    \mathbf{K}_a(\mathbf{k}) = \mathbf{B}_a^{\dagger}(\mathbf{k}) \mathbf{H}_{\mathrm{c}} \mathbf{B}_a(\mathbf{k}).
\end{equation}

The complete primitive-cell stiffness matrix is
\begin{equation}
    \boxed{ \mathbf{K}_{\mathrm{FCC}}(\mathbf{k}) = \sum_{a=1}^{6} \mathbf{B}_a^{\dagger}(\mathbf{k}) \mathbf{H}_{\mathrm{c}} \mathbf{B}_a(\mathbf{k}) }.
\end{equation}

Let $\mathbf{M}_{\mathrm{FCC}}^{(0)}$ denote the unstrained generalized inertial matrix of the primitive node--condensate system. The complete microscopic eigenvalue problem is
\begin{equation}
    \boxed{ \det \left[ \mathbf{K}_{\mathrm{FCC}}(\mathbf{k}) - \omega^2 \mathbf{M}_{\mathrm{FCC}}^{(0)} \right] = 0. }
\end{equation}

This equation contains the translational, microrotational, and internal-phase branches of the unstrained active FCC vacuum.

\subsection{Zero-Wave-Number Symmetry Conditions}

At
\begin{equation}
    \mathbf{k} = \mathbf{0},
\end{equation}
uniform translation of the entire lattice produces no contact deformation. The spectrum must therefore contain three exact translational null modes:
\begin{equation}
    \omega_1(\mathbf{0}) = \omega_2(\mathbf{0}) = \omega_3(\mathbf{0}) = 0.
\end{equation}

Because the phase contribution depends only on differences of the form
\begin{equation}
    \varphi_j - \varphi_i,
\end{equation}
a uniform internal-phase displacement also produces no phase-contact energy:
\begin{equation}
    \omega_{\varphi}(\mathbf{0}) = 0,
\end{equation}
unless a separate local phase-locking term is introduced.

A uniform independent microrotation is not necessarily a zero-energy mode. Frame indifference requires zero contact energy only when the nodal microrotation is accompanied by the compatible rigid rotation of the displacement field.

Mechanical stability of the discrete active phase requires
\begin{equation}
    \boxed{ \omega_s^2(\mathbf{k}) \geq 0 }
\end{equation}
for every branch $s$ and every wave vector $\mathbf{k}$ within the first Brillouin zone.

The quadratic stability conditions established for the individual contact tangent matrix are therefore necessary local conditions for stability of the assembled lattice.

\subsection{Long-Wavelength Affine Deformation}

The leading continuum tensor is obtained from a spatially affine displacement field,
\begin{equation}
    \mathbf{u}(\mathbf{x}) = \mathbf{H}\mathbf{x},
\end{equation}
where
\begin{equation}
    \mathbf{H} = \nabla\mathbf{u}.
\end{equation}

Decompose the displacement gradient into its symmetric and antisymmetric parts:
\begin{equation}
    \mathbf{H} = \boldsymbol{\varepsilon} + \mathbf{W},
\end{equation}
with
\begin{equation}
    \boldsymbol{\varepsilon} = \frac{1}{2} \left( \mathbf{H}+\mathbf{H}^{\mathsf{T}} \right),
\end{equation}
and
\begin{equation}
    \mathbf{W} = \frac{1}{2} \left( \mathbf{H}-\mathbf{H}^{\mathsf{T}} \right).
\end{equation}

The tensor $\boldsymbol{\varepsilon}$ is the infinitesimal strain, while $\mathbf{W}$ represents the macroscopic rigid rotation.

Let $\boldsymbol{\omega}$ denote the axial vector associated with $\mathbf{W}$, defined by
\begin{equation}
    \mathbf{W}\hat{\mathbf{n}} = \boldsymbol{\omega}\times\hat{\mathbf{n}}.
\end{equation}

Across contact $a$, the affine relative center displacement is
\begin{equation}
    \mathbf{u}(\mathbf{x}+\mathbf{d}_a) - \mathbf{u}(\mathbf{x}) = d_{\mathrm{nn}} \mathbf{H}\hat{\mathbf{n}}_a.
\end{equation}

For a spatially uniform nodal microrotation $\boldsymbol{\vartheta}$, the objective relative contact-point displacement becomes
\begin{equation}
    \Delta\mathbf{u}_a^{\mathrm{c}} = d_{\mathrm{nn}} \left[ \boldsymbol{\varepsilon}\hat{\mathbf{n}}_a + \left( \boldsymbol{\omega} - \boldsymbol{\vartheta} \right) \times \hat{\mathbf{n}}_a \right].
\end{equation}

The normal deformation is therefore
\begin{equation}
    \boxed{ \delta_{\mathrm{n}}^{(a)} = d_{\mathrm{nn}} \hat{\mathbf{n}}_a \cdot \boldsymbol{\varepsilon} \hat{\mathbf{n}}_a }.
\end{equation}

The tangential deformation is
\begin{equation}
    \boxed{ \boldsymbol{\delta}_{\mathrm{t}}^{(a)} = d_{\mathrm{nn}} \mathbf{P}_{\perp}^{(a)} \left[ \boldsymbol{\varepsilon}\hat{\mathbf{n}}_a + \left( \boldsymbol{\omega} - \boldsymbol{\vartheta} \right) \times \hat{\mathbf{n}}_a \right] },
\end{equation}
where
\begin{equation}
    \mathbf{P}_{\perp}^{(a)} = \mathbf{I} - \hat{\mathbf{n}}_a \hat{\mathbf{n}}_a^{\mathsf{T}}.
\end{equation}

\subsection{Relaxation of the Uniform Microrotation}

The six FCC contact directions satisfy the second-order identity
\begin{equation}
    \boxed{ \sum_{a=1}^{6} n_i^{(a)}n_j^{(a)} = 2\delta_{ij}. }
\end{equation}

For a symmetric strain tensor, this identity gives
\begin{equation}
    \sum_{a=1}^{6} \left( \boldsymbol{\varepsilon}\hat{\mathbf{n}}_a \right) \cdot \left[ \left( \boldsymbol{\omega} - \boldsymbol{\vartheta} \right) \times \hat{\mathbf{n}}_a \right] = 0.
\end{equation}

Consequently, the strain-dependent and relative-rotation-dependent parts of the summed tangential energy decouple at leading affine order.

The rotational contribution is positive and proportional to
\begin{equation}
    \left| \boldsymbol{\omega} - \boldsymbol{\vartheta} \right|^2.
\end{equation}

Minimization with respect to the uniform nodal microrotation therefore gives
\begin{equation}
    \boxed{ \boldsymbol{\vartheta} = \boldsymbol{\omega}. }
\end{equation}

The microrotation follows the macroscopic rigid rotation, removing rigid rotational motion from the contact strain.

After relaxation,
\begin{equation}
    \boxed{ \boldsymbol{\delta}_{\mathrm{t}}^{(a)} = d_{\mathrm{nn}} \mathbf{P}_{\perp}^{(a)} \boldsymbol{\varepsilon} \hat{\mathbf{n}}_a. }
\end{equation}

For a homogeneous affine deformation, neighboring cells possess the same microrotation and internal phase. Hence,
\begin{equation}
    \boldsymbol{\delta}_{\mathrm{b}}^{(a)} = \mathbf{0},
\end{equation}
\begin{equation}
    \delta_{\mathrm{tw}}^{(a)} = 0,
\end{equation}
and
\begin{equation}
    \Delta\varphi_a = 0.
\end{equation}

The bending, twisting, phase, and tangential--bending sectors therefore do not contribute to the leading classical elastic tensor. They remain active in rotational, torsional, phase, and finite-wave-number branches.

\subsection{Homogenized FCC Scaffold Energy}

The contact energy developed previously is defined per unit participating interface area. Let
\begin{equation}
    A_{\mathrm{ct}}
\end{equation}
denote the effective participating area assigned to one independent nearest-neighbor contact.

This is a local homogenization measure. Its value does not enter the isotropy condition derived below and is not assigned independently at this stage.

The solid-scaffold strain-energy density is
\begin{equation}
    W_{\mathrm{FCC}} = \frac{A_{\mathrm{ct}}}{V_{\mathrm{p}}} \sum_{a=1}^{6} \left[ \frac{1}{2} \kappa_{\mathrm{n}} \left( \delta_{\mathrm{n}}^{(a)} \right)^2 + \frac{1}{2} \kappa_{\mathrm{t}} \left| \boldsymbol{\delta}_{\mathrm{t}}^{(a)} \right|^2 \right].
\end{equation}

Using the relaxed affine contact kinematics,
\begin{equation}
    \left| \boldsymbol{\delta}_{\mathrm{t}}^{(a)} \right|^2 = d_{\mathrm{nn}}^2 \left[ \left| \boldsymbol{\varepsilon} \hat{\mathbf{n}}_a \right|^2 - \left( \hat{\mathbf{n}}_a \cdot \boldsymbol{\varepsilon} \hat{\mathbf{n}}_a \right)^2 \right],
\end{equation}
so the energy density becomes
\begin{equation}
\begin{aligned}
    W_{\mathrm{FCC}} ={}& \frac{ A_{\mathrm{ct}}d_{\mathrm{nn}}^2 }{ 2V_{\mathrm{p}} } \sum_{a=1}^{6} \Bigg[ \kappa_{\mathrm{t}} \left| \boldsymbol{\varepsilon} \hat{\mathbf{n}}_a \right|^2 \\
    &\hspace{29mm} + \left( \kappa_{\mathrm{n}} - \kappa_{\mathrm{t}} \right) \left( \hat{\mathbf{n}}_a \cdot \boldsymbol{\varepsilon} \hat{\mathbf{n}}_a \right)^2 \Bigg].
\end{aligned}
\end{equation}

The required FCC directional sums are
\begin{equation}
    \boxed{ \sum_{a=1}^{6} \left| \boldsymbol{\varepsilon} \hat{\mathbf{n}}_a \right|^2 = 2 \boldsymbol{\varepsilon} : \boldsymbol{\varepsilon}, }
\end{equation}
and
\begin{equation}
\begin{aligned}
    \sum_{a=1}^{6} \left( \hat{\mathbf{n}}_a \cdot \boldsymbol{\varepsilon} \hat{\mathbf{n}}_a \right)^2 ={}& \varepsilon_{xx}^2 + \varepsilon_{yy}^2 + \varepsilon_{zz}^2 \\
    &+ \varepsilon_{xx}\varepsilon_{yy} + \varepsilon_{xx}\varepsilon_{zz} + \varepsilon_{yy}\varepsilon_{zz} \\
    &+ 2 \left( \varepsilon_{xy}^2 + \varepsilon_{xz}^2 + \varepsilon_{yz}^2 \right).
\end{aligned}
\end{equation}

Substitution yields
\begin{equation}
\begin{aligned}
    W_{\mathrm{FCC}} = \frac{ A_{\mathrm{ct}}d_{\mathrm{nn}}^2 }{ 2V_{\mathrm{p}} } \Bigg\{ & \left( \kappa_{\mathrm{n}} + \kappa_{\mathrm{t}} \right) \left( \varepsilon_{xx}^2 + \varepsilon_{yy}^2 + \varepsilon_{zz}^2 \right) \\
    &+ \left( \kappa_{\mathrm{n}} - \kappa_{\mathrm{t}} \right) \left( \varepsilon_{xx}\varepsilon_{yy} + \varepsilon_{xx}\varepsilon_{zz} + \varepsilon_{yy}\varepsilon_{zz} \right) \\
    &+ 2 \left( \kappa_{\mathrm{n}} + \kappa_{\mathrm{t}} \right) \left( \varepsilon_{xy}^2 + \varepsilon_{xz}^2 + \varepsilon_{yz}^2 \right) \Bigg\}.
\end{aligned}
\end{equation}

\subsection{Emergent Cubic Elastic Constants}

The strain-energy density of a cubic continuum is
\begin{equation}
\begin{aligned}
    W_{\mathrm{cub}} ={}& \frac{1}{2} C_{11} \left( \varepsilon_{xx}^2 + \varepsilon_{yy}^2 + \varepsilon_{zz}^2 \right) \\
    &+ C_{12} \left( \varepsilon_{xx}\varepsilon_{yy} + \varepsilon_{xx}\varepsilon_{zz} + \varepsilon_{yy}\varepsilon_{zz} \right) \\
    &+ 2C_{44} \left( \varepsilon_{xy}^2 + \varepsilon_{xz}^2 + \varepsilon_{yz}^2 \right).
\end{aligned}
\end{equation}

Comparison with the homogenized FCC energy gives
\begin{equation}
    \boxed{ C_{11}^{(\mathrm{FCC})} = \frac{ A_{\mathrm{ct}}d_{\mathrm{nn}}^2 }{ V_{\mathrm{p}} } \left( \kappa_{\mathrm{n}} + \kappa_{\mathrm{t}} \right) },
\end{equation}
\begin{equation}
    \boxed{ C_{12}^{(\mathrm{FCC})} = \frac{ A_{\mathrm{ct}}d_{\mathrm{nn}}^2 }{ 2V_{\mathrm{p}} } \left( \kappa_{\mathrm{n}} - \kappa_{\mathrm{t}} \right) },
\end{equation}
and
\begin{equation}
    \boxed{ C_{44}^{(\mathrm{FCC})} = \frac{ A_{\mathrm{ct}}d_{\mathrm{nn}}^2 }{ 2V_{\mathrm{p}} } \left( \kappa_{\mathrm{n}} + \kappa_{\mathrm{t}} \right). }
\end{equation}

These are the leading elastic constants of the objective nearest-neighbor FCC scaffold after relaxation of the uniform microrotation.

\subsection{Mechanical Stability}

The leading cubic continuum is mechanically stable when
\begin{equation}
    C_{44} > 0,
\end{equation}
\begin{equation}
    C_{11} - C_{12} > 0,
\end{equation}
and
\begin{equation}
    C_{11} + 2C_{12} > 0.
\end{equation}

For the FCC scaffold,
\begin{equation}
    C_{44}^{(\mathrm{FCC})} = \frac{ A_{\mathrm{ct}}d_{\mathrm{nn}}^2 }{ 2V_{\mathrm{p}} } \left( \kappa_{\mathrm{n}} + \kappa_{\mathrm{t}} \right),
\end{equation}
\begin{equation}
    C_{11}^{(\mathrm{FCC})} - C_{12}^{(\mathrm{FCC})} = \frac{ A_{\mathrm{ct}}d_{\mathrm{nn}}^2 }{ 2V_{\mathrm{p}} } \left( \kappa_{\mathrm{n}} + 3\kappa_{\mathrm{t}} \right),
\end{equation}
and
\begin{equation}
    C_{11}^{(\mathrm{FCC})} + 2C_{12}^{(\mathrm{FCC})} = 2 \frac{ A_{\mathrm{ct}}d_{\mathrm{nn}}^2 }{ V_{\mathrm{p}} } \kappa_{\mathrm{n}}.
\end{equation}

Therefore,
\begin{equation}
    \boxed{ \kappa_{\mathrm{n}} > 0, \qquad \kappa_{\mathrm{t}} > 0 }
\end{equation}
are sufficient for stability of the leading affine FCC continuum.

\subsection{Condition for Macroscopic Isotropy}

The elastic anisotropy factor of a cubic solid is
\begin{equation}
    \boxed{ \mathcal{A}_{\mathrm{cub}} = \frac{ 2C_{44} }{ C_{11}-C_{12} }. }
\end{equation}

For the homogenized FCC scaffold,
\begin{equation}
    \boxed{ \mathcal{A}_{\mathrm{cub}} = \frac{ 2 \left( \kappa_{\mathrm{n}} + \kappa_{\mathrm{t}} \right) }{ \kappa_{\mathrm{n}} + 3\kappa_{\mathrm{t}} }. }
\end{equation}

Macroscopic isotropy requires
\begin{equation}
    \mathcal{A}_{\mathrm{cub}} = 1.
\end{equation}

Consequently,
\begin{equation}
    2 \left( \kappa_{\mathrm{n}} + \kappa_{\mathrm{t}} \right) = \kappa_{\mathrm{n}} + 3\kappa_{\mathrm{t}},
\end{equation}
and therefore
\begin{equation}
    \boxed{ \kappa_{\mathrm{n}} = \kappa_{\mathrm{t}}. }
\end{equation}

The active FCC scaffold is thus isotropic at leading continuum order precisely when its incremental normal and tangential contact stiffnesses are equal.

Define the common contact coefficient by
\begin{equation}
    \boxed{ \kappa_{\mathrm{c}} \equiv \kappa_{\mathrm{n}} = \kappa_{\mathrm{t}}. }
\end{equation}

Under this condition,
\begin{equation}
    C_{11}^{(\mathrm{FCC})} = 2 \frac{ A_{\mathrm{ct}}d_{\mathrm{nn}}^2 }{ V_{\mathrm{p}} } \kappa_{\mathrm{c}},
\end{equation}
\begin{equation}
    C_{12}^{(\mathrm{FCC})} = 0,
\end{equation}
and
\begin{equation}
    C_{44}^{(\mathrm{FCC})} = \frac{ A_{\mathrm{ct}}d_{\mathrm{nn}}^2 }{ V_{\mathrm{p}} } \kappa_{\mathrm{c}}.
\end{equation}

Define the homogenized scaffold shear modulus by
\begin{equation}
    \boxed{ K_{\mathrm{sh}} = \frac{ A_{\mathrm{ct}}d_{\mathrm{nn}}^2 }{ V_{\mathrm{p}} } \kappa_{\mathrm{c}}. }
\end{equation}

The isotropic scaffold constants then become
\begin{equation}
    C_{11}^{(\mathrm{FCC})} = 2K_{\mathrm{sh}},
\end{equation}
\begin{equation}
    C_{12}^{(\mathrm{FCC})} = 0,
\end{equation}
and
\begin{equation}
    C_{44}^{(\mathrm{FCC})} = K_{\mathrm{sh}}.
\end{equation}

The nearest-neighbor solid scaffold therefore supplies isotropic shear rigidity but does not, by itself, produce a nonzero first Lam\'{e} coefficient.

\subsection{Interstitial-Condensate Volumetric Response}

The active vacuum is a dual-phase medium. The solid FCC scaffold supplies transverse contact rigidity, while the interstitial condensate supplies the complementary volumetric and inertial response.

Let
\begin{equation}
    B_{\mathrm{IC}}
\end{equation}
denote the leading isotropic volumetric coefficient of the interstitial-condensate sector. Its affine elastic energy density is
\begin{equation}
    \boxed{ W_{\mathrm{IC}} = \frac{1}{2} B_{\mathrm{IC}} \left( \operatorname{tr} \boldsymbol{\varepsilon} \right)^2. }
\end{equation}

This contribution adds $B_{\mathrm{IC}}$ to both $C_{11}$ and $C_{12}$, while leaving $C_{44}$ unchanged:
\begin{equation}
    C_{11} = C_{11}^{(\mathrm{FCC})} + B_{\mathrm{IC}},
\end{equation}
\begin{equation}
    C_{12} = C_{12}^{(\mathrm{FCC})} + B_{\mathrm{IC}},
\end{equation}
and
\begin{equation}
    C_{44} = C_{44}^{(\mathrm{FCC})}.
\end{equation}

Because the same volumetric coefficient is added to $C_{11}$ and $C_{12}$,
\begin{equation}
    C_{11}-C_{12}
\end{equation}
is unchanged. The condensate therefore does not alter the isotropy condition obtained from the FCC scaffold.

Under
\begin{equation}
    \kappa_{\mathrm{n}} = \kappa_{\mathrm{t}},
\end{equation}
the complete active-vacuum elastic constants are
\begin{equation}
    \boxed{ C_{11} = B_{\mathrm{IC}} + 2K_{\mathrm{sh}}, }
\end{equation}
\begin{equation}
    \boxed{ C_{12} = B_{\mathrm{IC}}, }
\end{equation}
and
\begin{equation}
    \boxed{ C_{44} = K_{\mathrm{sh}}. }
\end{equation}

The leading active-vacuum elastic tensor is consequently
\begin{equation}
    \boxed{ C_{ijkl} = B_{\mathrm{IC}} \delta_{ij}\delta_{kl} + K_{\mathrm{sh}} \left( \delta_{ik}\delta_{jl} + \delta_{il}\delta_{jk} \right). }
\end{equation}

The corresponding total bulk modulus is
\begin{equation}
    \boxed{ K_{\mathrm{bulk}} = B_{\mathrm{IC}} + \frac{2}{3}K_{\mathrm{sh}}. }
\end{equation}

The dual-phase constitutive division is therefore
\begin{equation}
    \boxed{ \text{FCC nodal scaffold} \longrightarrow \text{shear rigidity}, }
\end{equation}
and
\begin{equation}
    \boxed{ \text{interstitial condensate} \longrightarrow \text{volumetric response and inertia}. }
\end{equation}

\subsection{Long-Wavelength Acoustic Tensor}

For propagation in the unit direction
\begin{equation}
    \hat{\mathbf{s}} = \frac{\mathbf{k}}{|\mathbf{k}|},
\end{equation}
the acoustic tensor is
\begin{equation}
    \Gamma_{ij} = C_{ikjl}s_k s_l.
\end{equation}

Using the isotropic active-vacuum tensor gives
\begin{equation}
    \boxed{ \Gamma_{ij} = K_{\mathrm{sh}}\delta_{ij} + \left( B_{\mathrm{IC}} + K_{\mathrm{sh}} \right) s_i s_j. }
\end{equation}

Let
\begin{equation}
    \rho_{\mathrm{eff}}^{(0)}
\end{equation}
denote the unstrained effective translational inertia density of the homogenized node--condensate system.

For a transverse polarization satisfying
\begin{equation}
    \hat{\mathbf{s}}\cdot\mathbf{e}_T = 0,
\end{equation}
the acoustic equation gives
\begin{equation}
    \omega_T^2 = \frac{ K_{\mathrm{sh}} }{ \rho_{\mathrm{eff}}^{(0)} } k^2 + \mathcal{O}(k^4).
\end{equation}

The unstrained transverse speed is therefore
\begin{equation}
    \boxed{ c_T^{(0)} = \sqrt{ \frac{ K_{\mathrm{sh}} }{ \rho_{\mathrm{eff}}^{(0)} } }. }
\end{equation}

Consistency with the established active-vacuum elastic capacity requires
\begin{equation}
    c_T^{(0)} = c_{\mathrm{E}}.
\end{equation}

Hence,
\begin{equation}
    \boxed{ K_{\mathrm{sh}} = \rho_{\mathrm{eff}}^{(0)} c_{\mathrm{E}}^2. }
\end{equation}

In terms of the contact parameters,
\begin{equation}
    \boxed{ \frac{ A_{\mathrm{ct}}d_{\mathrm{nn}}^2 }{ V_{\mathrm{p}} } \kappa_{\mathrm{c}} = \rho_{\mathrm{eff}}^{(0)} c_{\mathrm{E}}^2. }
\end{equation}

The isotropy condition fixes the ratio of normal to tangential contact stiffness, while the transverse-wave normalization fixes the corresponding stiffness--inertia combination.

The longitudinal branch satisfies
\begin{equation}
    \omega_L^2 = \frac{ B_{\mathrm{IC}} + 2K_{\mathrm{sh}} }{ \rho_{\mathrm{eff}}^{(0)} } k^2 + \mathcal{O}(k^4).
\end{equation}

Its speed depends additionally on the condensate volumetric coefficient and is not required to equal the transverse elastic limit $c_{\mathrm{E}}$.

\subsection{Gravitationally Constrained Dynamic Tensor}

The endpoint-transfer Jacobian for one nearest-neighbor contact is
\begin{equation}
    \mathcal{C}_{ij} = \left[ 1-\alpha_{\mathrm{G}}^{(i)} \right] \left[ 1-\alpha_{\mathrm{G}}^{(j)} \right].
\end{equation}

For a gravitational field varying negligibly over one nearest-neighbor separation,
\begin{equation}
    \mathcal{C} = \left( 1-\alpha_{\mathrm{G}} \right)^2.
\end{equation}

The reduced translational inertia density is therefore
\begin{equation}
    \boxed{ \rho_{\mathrm{eff}} \left( \alpha_{\mathrm{G}} \right) = \frac{ \rho_{\mathrm{eff}}^{(0)} }{ \left( 1-\alpha_{\mathrm{G}} \right)^4 }. }
\end{equation}

At the retained active-phase constitutive order, the contact Hessian and the homogenized elastic tensor remain unchanged:
\begin{equation}
    \boxed{ C_{ijkl} \left( \alpha_{\mathrm{G}} \right) = C_{ijkl}(0), \qquad 0 \leq \alpha_{\mathrm{G}} < 1. }
\end{equation}

The gravitationally loaded acoustic equation is
\begin{equation}
    \Gamma_{ij}e_j = \rho_{\mathrm{eff}} \left( \alpha_{\mathrm{G}} \right) c^2e_i.
\end{equation}

For either transverse polarization,
\begin{equation}
    c_T^2 \left( \alpha_{\mathrm{G}} \right) = \frac{ K_{\mathrm{sh}} }{ \rho_{\mathrm{eff}} \left( \alpha_{\mathrm{G}} \right) }.
\end{equation}

Using
\begin{equation}
    K_{\mathrm{sh}} = \rho_{\mathrm{eff}}^{(0)} c_{\mathrm{E}}^2,
\end{equation}
one obtains
\begin{equation}
    c_T^2 \left( \alpha_{\mathrm{G}} \right) = c_{\mathrm{E}}^2 \left( 1-\alpha_{\mathrm{G}} \right)^4.
\end{equation}

Therefore,
\begin{equation}
    \boxed{ c_T \left( \alpha_{\mathrm{G}} \right) = c_{\mathrm{E}} \left( 1-\alpha_{\mathrm{G}} \right)^2. }
\end{equation}

Because the homogenized tensor is isotropic, this result applies equally to both transverse polarizations and to every long-wavelength propagation direction.

The one-contact endpoint-transfer law therefore survives the full twelve-neighbor FCC homogenization.

\subsection{Higher-Order Contact Sectors}

The leading affine elastic tensor depends only on
\begin{equation}
    \kappa_{\mathrm{n}}
\end{equation}
and
\begin{equation}
    \kappa_{\mathrm{t}}.
\end{equation}

The remaining contact coefficients,
\begin{equation}
    \kappa_{\mathrm{b}}, \qquad \kappa_{\mathrm{tw}}, \qquad \kappa_{\mathrm{tb}}, \qquad \kappa_{\varphi},
\end{equation}
enter when the microrotation and phase vary between neighboring cells.

Their contributions begin in:
\begin{itemize}
    \item rotational and torsional lattice branches;
    \item phase-propagation branches;
    \item finite-wave-number dispersion;
    \item strain-gradient elasticity;
    \item localized topological-defect structure;
    \item chiral response when parity-odd terms are admitted.
\end{itemize}

These sectors do not alter the leading affine tensor derived above because
\begin{equation}
    \boldsymbol{\delta}_{\mathrm{b}}^{(a)} = \mathbf{0}, \qquad \delta_{\mathrm{tw}}^{(a)} = 0, \qquad \Delta\varphi_a = 0
\end{equation}
for a homogeneous long-wavelength deformation after microrotational relaxation.

They remain essential to the complete non-affine microphysics of the active vacuum.

\subsection{Constitutive Result}

The twelve-neighbor FCC homogenization produces the scaffold elastic constants
\begin{equation}
    \boxed{ C_{11}^{(\mathrm{FCC})} = \frac{ A_{\mathrm{ct}}d_{\mathrm{nn}}^2 }{ V_{\mathrm{p}} } \left( \kappa_{\mathrm{n}} + \kappa_{\mathrm{t}} \right), }
\end{equation}
\begin{equation}
    \boxed{ C_{12}^{(\mathrm{FCC})} = \frac{ A_{\mathrm{ct}}d_{\mathrm{nn}}^2 }{ 2V_{\mathrm{p}} } \left( \kappa_{\mathrm{n}} - \kappa_{\mathrm{t}} \right), }
\end{equation}
and
\begin{equation}
    \boxed{ C_{44}^{(\mathrm{FCC})} = \frac{ A_{\mathrm{ct}}d_{\mathrm{nn}}^2 }{ 2V_{\mathrm{p}} } \left( \kappa_{\mathrm{n}} + \kappa_{\mathrm{t}} \right). }
\end{equation}

The cubic anisotropy factor is
\begin{equation}
    \boxed{ \mathcal{A}_{\mathrm{cub}} = \frac{ 2 \left( \kappa_{\mathrm{n}} + \kappa_{\mathrm{t}} \right) }{ \kappa_{\mathrm{n}} + 3\kappa_{\mathrm{t}} }. }
\end{equation}

The leading continuum is isotropic if and only if
\begin{equation}
    \boxed{ \kappa_{\mathrm{n}} = \kappa_{\mathrm{t}}. }
\end{equation}

Under that condition, the FCC scaffold supplies the shear modulus
\begin{equation}
    \boxed{ K_{\mathrm{sh}} = \frac{ A_{\mathrm{ct}}d_{\mathrm{nn}}^2 }{ V_{\mathrm{p}} } \kappa_{\mathrm{c}}, }
\end{equation}
while the interstitial condensate supplies the complementary volumetric coefficient $B_{\mathrm{IC}}$.

The complete leading elastic tensor is
\begin{equation}
    \boxed{ C_{ijkl} = B_{\mathrm{IC}} \delta_{ij}\delta_{kl} + K_{\mathrm{sh}} \left( \delta_{ik}\delta_{jl} + \delta_{il}\delta_{jk} \right). }
\end{equation}

The unstrained transverse normalization is
\begin{equation}
    \boxed{ K_{\mathrm{sh}} = \rho_{\mathrm{eff}}^{(0)} c_{\mathrm{E}}^2. }
\end{equation}

Gravitational loading modifies the effective inertia according to
\begin{equation}
    \boxed{ \rho_{\mathrm{eff}} \left( \alpha_{\mathrm{G}} \right) = \frac{ \rho_{\mathrm{eff}}^{(0)} }{ \left( 1-\alpha_{\mathrm{G}} \right)^4 }, }
\end{equation}
while leaving the active-phase elastic tensor unchanged at leading order.

The resulting three-dimensional transverse propagation law is
\begin{equation}
    \boxed{ c_P = c_{\mathrm{E}} \left( 1-\alpha_{\mathrm{G}} \right)^2. }
\end{equation}

The active-vacuum microphysical chain is therefore
\begin{equation}
    \boxed{
    \begin{array}{c}
    \text{objective FCC contact law} \\
    \longrightarrow \text{twelve-neighbor homogenization} \\
    \longrightarrow \text{isotropy condition} \\
    \longrightarrow \text{dual-phase elastic tensor} \\
    \longrightarrow \text{gravitationally constrained inertia} \\
    \longrightarrow c_P.
    \end{array}
    }
\end{equation}

The absolute participating contact area, the condensate volumetric coefficient, and the higher-order rotational and phase coefficients remain to be fixed by the geometric and energetic calibration of the complete active-vacuum contact cell. The tensor structure, isotropy condition, and gravitational propagation scaling do not depend on that remaining absolute calibration.

\section{Geometric Calibration of the FCC Contact Cell}

\subsection{Scope}

The preceding homogenization expressed the leading active-vacuum shear modulus in terms of the effective participating contact area $A_{\mathrm{ct}}$ and the isotropic contact coefficient $\kappa_{\mathrm{c}}$:
\begin{equation}
    K_{\mathrm{sh}} = \frac{ A_{\mathrm{ct}}d_{\mathrm{nn}}^2 }{ V_{\mathrm{p}} } \kappa_{\mathrm{c}},
\end{equation}
where
\begin{equation}
    d_{\mathrm{nn}} = 2r_{\mathrm{D}}
\end{equation}
is the FCC nearest-neighbor separation and
\begin{equation}
    V_{\mathrm{p}} = 4\sqrt{2}r_{\mathrm{D}}^3
\end{equation}
is the primitive-cell volume.

The directional transverse contact law established
\begin{equation}
    \kappa_{\mathrm{c}} = \frac{\tau_{\mathrm{D}}}{d_{\mathrm{nn}}^2},
\end{equation}
under the isotropy condition
\begin{equation}
    \kappa_{\mathrm{n}} = \kappa_{\mathrm{t}} = \kappa_{\mathrm{c}}.
\end{equation}

The remaining quantity $A_{\mathrm{ct}}$ cannot be identified with the entire geometric Wigner--Seitz face without first enforcing consistency with the independently derived inertia of one Soft Floor node. The purpose of the present section is therefore to determine the effective contact area by matching the surface-inertia description of the active interface to the localized node mass $M_{\mathrm{D}}$.

This calibration removes $A_{\mathrm{ct}}$, $\kappa_{\mathrm{c}}$, and the unstrained acoustic inertia density from the transverse contact theory.

\subsection{Wigner--Seitz Geometry of the FCC Node}

The Wigner--Seitz cell of an FCC lattice is a rhombic dodecahedron. Its twelve faces correspond one-to-one with the twelve nearest-neighbor directions.

For nearest-neighbor separation $d_{\mathrm{nn}}$, each rhombic face has diagonals
\begin{equation}
    d_{\mathrm{nn}}
\end{equation}
and
\begin{equation}
    \frac{d_{\mathrm{nn}}}{\sqrt{2}}.
\end{equation}

The area of one Wigner--Seitz face is therefore
\begin{equation}
    A_{\mathrm{WS}} = \frac{1}{2} d_{\mathrm{nn}} \frac{d_{\mathrm{nn}}}{\sqrt{2}}.
\end{equation}

Hence,
\begin{equation}
    \boxed{ A_{\mathrm{WS}} = \frac{d_{\mathrm{nn}}^2}{2\sqrt{2}} }.
\end{equation}

Using
\begin{equation}
    d_{\mathrm{nn}} = 2r_{\mathrm{D}},
\end{equation}
one obtains
\begin{equation}
    \boxed{ A_{\mathrm{WS}} = \sqrt{2}r_{\mathrm{D}}^2 }.
\end{equation}

The ratio of one full Wigner--Seitz face to the primitive-cell volume is
\begin{equation}
    \frac{A_{\mathrm{WS}}}{V_{\mathrm{p}}} = \frac{ \sqrt{2}r_{\mathrm{D}}^2 }{ 4\sqrt{2}r_{\mathrm{D}}^3 }.
\end{equation}

Therefore,
\begin{equation}
    \boxed{ \frac{A_{\mathrm{WS}}}{V_{\mathrm{p}}} = \frac{1}{4r_{\mathrm{D}}} }.
\end{equation}

This is the geometric surface-to-volume ratio associated with one complete nearest-neighbor face. It provides an upper geometric area for one directional transfer channel, but it does not by itself determine how much of that face participates in the reduced transverse state-transfer coordinate.

\subsection{Inertial Definition of the Participating Contact Area}

The surface inertia $\mu_{\mathrm{D}}$ is the kinetic coefficient per unit active interface area. If the effective area participating in one primitive translational transfer is $A_{\mathrm{ct}}$, the corresponding integrated inertial coefficient is
\begin{equation}
    m_{\mathrm{ct}} = \mu_{\mathrm{D}}A_{\mathrm{ct}}.
\end{equation}

The primitive translational coordinate describes the motion of one active Soft Floor node. Its integrated inertia must therefore equal the independently derived node mass $M_{\mathrm{D}}$:
\begin{equation}
    \boxed{ m_{\mathrm{ct}} = M_{\mathrm{D}} }.
\end{equation}

Consequently,
\begin{equation}
    \mu_{\mathrm{D}}A_{\mathrm{ct}} = M_{\mathrm{D}},
\end{equation}
and the effective participating contact area is fixed by
\begin{equation}
    \boxed{ A_{\mathrm{ct}} = \frac{M_{\mathrm{D}}}{\mu_{\mathrm{D}}} }.
\end{equation}

This relation is an inertial matching condition. It does not identify $A_{\mathrm{ct}}$ with the literal geometric area of the spherical contact point. Rather, $A_{\mathrm{ct}}$ is the effective interface-area measure multiplying $\mu_{\mathrm{D}}$ in the reduced state-transfer action.

\subsection{Reduction to the Substrate Area Scale}

Appendix C gives the node mass as
\begin{equation}
    M_{\mathrm{D}} = \frac{h}{c_{\mathrm{E}}r_{\mathrm{D}}} = \frac{2\pi\hbar}{c_{\mathrm{E}}r_{\mathrm{D}}},
\end{equation}
and the surface inertia as
\begin{equation}
    \mu_{\mathrm{D}} = \frac{c_{\mathrm{E}}^2}{Gr_{\mathrm{D}}}.
\end{equation}

Substitution into the inertial matching condition gives
\begin{equation}
    A_{\mathrm{ct}} = \frac{ 2\pi\hbar/(c_{\mathrm{E}}r_{\mathrm{D}}) }{ c_{\mathrm{E}}^2/(Gr_{\mathrm{D}}) }.
\end{equation}

The Soft Floor radius cancels:
\begin{equation}
    A_{\mathrm{ct}} = \frac{ 2\pi\hbar G }{ c_{\mathrm{E}}^3 }.
\end{equation}

The substrate lattice scale satisfies
\begin{equation}
    \ell_{\mathrm{S}}^2 = \frac{\hbar G}{c_{\mathrm{E}}^3}.
\end{equation}

Therefore,
\begin{equation}
    \boxed{ A_{\mathrm{ct}} = 2\pi\ell_{\mathrm{S}}^2 }.
\end{equation}

The effective area governing one FCC state-transfer coordinate is thus fixed by the substrate area scale rather than by the full Soft Floor face area.

The equivalent circular radius of the contact measure is
\begin{equation}
    r_{\mathrm{ct}} = \sqrt{ \frac{A_{\mathrm{ct}}}{\pi} }.
\end{equation}

Hence,
\begin{equation}
    \boxed{ r_{\mathrm{ct}} = \sqrt{2}\ell_{\mathrm{S}} }.
\end{equation}

Using the calibrated substrate scale, this corresponds approximately to
\begin{equation}
    A_{\mathrm{ct}} \approx 1.64\times 10^{-69} \ \mathrm{m}^2,
\end{equation}
and
\begin{equation}
    r_{\mathrm{ct}} \approx 2.29\times 10^{-35} \ \mathrm{m}.
\end{equation}

The contact-cell calibration therefore identifies the active transfer footprint with the substrate-scale anchoring of the Soft Floor node.

\subsection{Fraction of the Wigner--Seitz Face Participating in Transfer}

Define the dimensionless face-participation fraction by
\begin{equation}
    f_{\mathrm{ct}} = \frac{A_{\mathrm{ct}}}{A_{\mathrm{WS}}}.
\end{equation}

Using
\begin{equation}
    A_{\mathrm{ct}} = 2\pi\ell_{\mathrm{S}}^2
\end{equation}
and
\begin{equation}
    A_{\mathrm{WS}} = \sqrt{2}r_{\mathrm{D}}^2,
\end{equation}
one obtains
\begin{equation}
    f_{\mathrm{ct}} = \frac{ 2\pi\ell_{\mathrm{S}}^2 }{ \sqrt{2}r_{\mathrm{D}}^2 }.
\end{equation}

Therefore,
\begin{equation}
    \boxed{ f_{\mathrm{ct}} = \sqrt{2}\pi \left( \frac{\ell_{\mathrm{S}}}{r_{\mathrm{D}}} \right)^2 }.
\end{equation}

Because
\begin{equation}
    \ell_{\mathrm{S}} \ll r_{\mathrm{D}},
\end{equation}
the active transfer area occupies only a small fraction of the geometric Wigner--Seitz face:
\begin{equation}
    \boxed{ f_{\mathrm{ct}} \ll 1 }.
\end{equation}

For the calibrated TFT scales,
\begin{equation}
    f_{\mathrm{ct}} \approx 2.09\times 10^{-17}.
\end{equation}

Thus, the complete Wigner--Seitz face defines the nearest-neighbor geometric boundary, whereas the kinetic transfer coordinate is concentrated into a substrate-scale effective channel within that boundary.

Assigning the entire Wigner--Seitz face to the primitive translational inertia would instead give
\begin{equation}
    \mu_{\mathrm{D}}A_{\mathrm{WS}} = \frac{A_{\mathrm{WS}}}{A_{\mathrm{ct}}} M_{\mathrm{D}},
\end{equation}
which would exceed the independently derived node inertia by the factor
\begin{equation}
    \frac{A_{\mathrm{WS}}}{A_{\mathrm{ct}}} = \frac{1}{f_{\mathrm{ct}}}.
\end{equation}

The full face therefore cannot serve as the inertial area of one reduced transfer coordinate.

\subsection{Integrated Contact Stiffness}

The isotropic contact coefficient is
\begin{equation}
    \kappa_{\mathrm{c}} = \frac{\tau_{\mathrm{D}}}{d_{\mathrm{nn}}^2}.
\end{equation}

Its dimensions are
\begin{equation}
    [\kappa_{\mathrm{c}}] = \mathrm{N/m^3}.
\end{equation}

The integrated stiffness of one effective nearest-neighbor transfer channel is
\begin{equation}
    k_{\mathrm{ct}} = A_{\mathrm{ct}}\kappa_{\mathrm{c}}.
\end{equation}

Therefore,
\begin{equation}
    k_{\mathrm{ct}} = A_{\mathrm{ct}} \frac{\tau_{\mathrm{D}}}{d_{\mathrm{nn}}^2}.
\end{equation}

Using
\begin{equation}
    A_{\mathrm{ct}} = 2\pi\ell_{\mathrm{S}}^2
\end{equation}
and
\begin{equation}
    d_{\mathrm{nn}} = 2r_{\mathrm{D}},
\end{equation}
gives
\begin{equation}
    k_{\mathrm{ct}} = 2\pi\ell_{\mathrm{S}}^2 \frac{\tau_{\mathrm{D}}}{4r_{\mathrm{D}}^2}.
\end{equation}

Hence,
\begin{equation}
    \boxed{ k_{\mathrm{ct}} = \frac{\pi}{2} \tau_{\mathrm{D}} \left( \frac{\ell_{\mathrm{S}}}{r_{\mathrm{D}}} \right)^2 }.
\end{equation}

The discrete transverse contact action may consequently be written using the fully calibrated pair
\begin{equation}
    \boxed{ m_{\mathrm{ct}} = M_{\mathrm{D}} },
\end{equation}
and
\begin{equation}
    \boxed{ k_{\mathrm{ct}} = \frac{\pi}{2} \tau_{\mathrm{D}} \left( \frac{\ell_{\mathrm{S}}}{r_{\mathrm{D}}} \right)^2 }.
\end{equation}

No independent contact area or discrete spring coefficient remains.

\subsection{Discrete Propagation Normalization}

The long-wavelength speed of a nearest-neighbor contact sequence is
\begin{equation}
    c_T^{(0)} = d_{\mathrm{nn}} \sqrt{ \frac{k_{\mathrm{ct}}}{m_{\mathrm{ct}}} }.
\end{equation}

Using
\begin{equation}
    k_{\mathrm{ct}} = A_{\mathrm{ct}}\kappa_{\mathrm{c}},
\end{equation}
\begin{equation}
    m_{\mathrm{ct}} = A_{\mathrm{ct}}\mu_{\mathrm{D}},
\end{equation}
and
\begin{equation}
    \kappa_{\mathrm{c}} = \frac{\tau_{\mathrm{D}}}{d_{\mathrm{nn}}^2},
\end{equation}
one obtains
\begin{equation}
\begin{aligned}
    \left( c_T^{(0)} \right)^2 &= d_{\mathrm{nn}}^2 \frac{ A_{\mathrm{ct}}\kappa_{\mathrm{c}} }{ A_{\mathrm{ct}}\mu_{\mathrm{D}} } \\[3pt]
    &= d_{\mathrm{nn}}^2 \frac{ \tau_{\mathrm{D}}/d_{\mathrm{nn}}^2 }{ \mu_{\mathrm{D}} }.
\end{aligned}
\end{equation}

The contact area and lattice spacing cancel:
\begin{equation}
    \left( c_T^{(0)} \right)^2 = \frac{\tau_{\mathrm{D}}}{\mu_{\mathrm{D}}}.
\end{equation}

Using the established Soft Floor identity,
\begin{equation}
    \frac{\tau_{\mathrm{D}}}{\mu_{\mathrm{D}}} = c_{\mathrm{E}}^2,
\end{equation}
gives
\begin{equation}
    \boxed{ c_T^{(0)} = c_{\mathrm{E}} }.
\end{equation}

The invariant transverse propagation speed therefore follows from the calibrated contact mass and stiffness without an additional wave-speed fit.

\subsection{Primitive-Cell Acoustic Inertia Density}

The unstrained translational inertia density of the primitive FCC cell is
\begin{equation}
    \rho_{\mathrm{eff}}^{(0)} = \frac{M_{\mathrm{D}}}{V_{\mathrm{p}}}.
\end{equation}

Using
\begin{equation}
    M_{\mathrm{D}} = \rho_{\mathrm{D}} \left( \frac{4}{3}\pi r_{\mathrm{D}}^3 \right),
\end{equation}
and
\begin{equation}
    V_{\mathrm{p}} = 4\sqrt{2}r_{\mathrm{D}}^3,
\end{equation}
one obtains
\begin{equation}
    \rho_{\mathrm{eff}}^{(0)} = \rho_{\mathrm{D}} \frac{ (4/3)\pi r_{\mathrm{D}}^3 }{ 4\sqrt{2}r_{\mathrm{D}}^3 }.
\end{equation}

Therefore,
\begin{equation}
    \boxed{ \rho_{\mathrm{eff}}^{(0)} = \frac{\pi}{3\sqrt{2}} \rho_{\mathrm{D}} }.
\end{equation}

The FCC close-packing fraction is
\begin{equation}
    \beta_{\mathrm{geo}} = \frac{\pi}{\sqrt{18}} = \frac{\pi}{3\sqrt{2}}.
\end{equation}

Hence,
\begin{equation}
    \boxed{ \rho_{\mathrm{eff}}^{(0)} = \beta_{\mathrm{geo}}\rho_{\mathrm{D}} }.
\end{equation}

The homogenized translational inertia is therefore exactly the operational node density weighted by the solid FCC packing fraction.

This result connects the primitive-cell acoustic mass directly to the geometric partition of the active vacuum.

Using the action representation,
\begin{equation}
    M_{\mathrm{D}} = \frac{2\pi\hbar}{c_{\mathrm{E}}r_{\mathrm{D}}},
\end{equation}
the same density may be written as
\begin{equation}
    \boxed{ \rho_{\mathrm{eff}}^{(0)} = \frac{\pi\hbar}{2\sqrt{2}c_{\mathrm{E}}r_{\mathrm{D}}^4} }.
\end{equation}

\subsection{Absolute Homogenized Shear Modulus}

The homogenized scaffold shear modulus obtained previously is
\begin{equation}
    K_{\mathrm{sh}} = \frac{ A_{\mathrm{ct}}d_{\mathrm{nn}}^2 }{ V_{\mathrm{p}} } \kappa_{\mathrm{c}}.
\end{equation}

Using
\begin{equation}
    d_{\mathrm{nn}}^2\kappa_{\mathrm{c}} = \tau_{\mathrm{D}},
\end{equation}
this reduces to
\begin{equation}
    K_{\mathrm{sh}} = \frac{A_{\mathrm{ct}}}{V_{\mathrm{p}}} \tau_{\mathrm{D}}.
\end{equation}

Substitution of the calibrated contact area gives
\begin{equation}
    \boxed{ K_{\mathrm{sh}} = \frac{ 2\pi\ell_{\mathrm{S}}^2 }{ 4\sqrt{2}r_{\mathrm{D}}^3 } \tau_{\mathrm{D}} }.
\end{equation}

Therefore,
\begin{equation}
    \boxed{ K_{\mathrm{sh}} = \frac{\pi}{2\sqrt{2}} \frac{\ell_{\mathrm{S}}^2}{r_{\mathrm{D}}^3} \tau_{\mathrm{D}} }.
\end{equation}

Because
\begin{equation}
    \frac{\tau_{\mathrm{D}}}{\mu_{\mathrm{D}}} = c_{\mathrm{E}}^2,
\end{equation}
and
\begin{equation}
    \frac{ A_{\mathrm{ct}}\mu_{\mathrm{D}} }{ V_{\mathrm{p}} } = \frac{M_{\mathrm{D}}}{V_{\mathrm{p}}} = \rho_{\mathrm{eff}}^{(0)},
\end{equation}
the shear modulus satisfies
\begin{equation}
    \boxed{ K_{\mathrm{sh}} = \rho_{\mathrm{eff}}^{(0)} c_{\mathrm{E}}^2 }.
\end{equation}

Using the FCC packing result,
\begin{equation}
    \boxed{ K_{\mathrm{sh}} = \beta_{\mathrm{geo}} \rho_{\mathrm{D}} c_{\mathrm{E}}^2 }.
\end{equation}

An equivalent expression in terms of the action scale is
\begin{equation}
    \boxed{ K_{\mathrm{sh}} = \frac{ \pi\hbar c_{\mathrm{E}} }{ 2\sqrt{2}r_{\mathrm{D}}^4 } }.
\end{equation}

The absolute small-strain shear modulus is thus fixed entirely by the established substrate scale, Soft Floor radius, and action normalization.

\subsection{Distinction from the Gravitational Yield Tension}

The manuscript defines the volumetric gravitational yield tension by
\begin{equation}
    \tau_{\mathrm{D,bulk}} = \frac{\tau_{\mathrm{D}}}{r_{\mathrm{D}}}.
\end{equation}

The calibrated small-strain shear modulus is not equal to this yield scale.

Their ratio is
\begin{equation}
    \frac{K_{\mathrm{sh}}}{\tau_{\mathrm{D,bulk}}} = \frac{ A_{\mathrm{ct}}\tau_{\mathrm{D}}/V_{\mathrm{p}} }{ \tau_{\mathrm{D}}/r_{\mathrm{D}} }.
\end{equation}

Therefore,
\begin{equation}
    \frac{K_{\mathrm{sh}}}{\tau_{\mathrm{D,bulk}}} = \frac{ A_{\mathrm{ct}}r_{\mathrm{D}} }{ V_{\mathrm{p}} }.
\end{equation}

Using the calibrated geometry gives
\begin{equation}
    \boxed{ \frac{K_{\mathrm{sh}}}{\tau_{\mathrm{D,bulk}}} = \frac{\pi}{2\sqrt{2}} \left( \frac{\ell_{\mathrm{S}}}{r_{\mathrm{D}}} \right)^2 }.
\end{equation}

Thus,
\begin{equation}
    \boxed{ K_{\mathrm{sh}} \ll \tau_{\mathrm{D,bulk}} }.
\end{equation}

The difference is mechanical rather than dimensional:
\begin{itemize}
    \item $K_{\mathrm{sh}}$ is the tangent modulus governing small transverse reconfiguration through the substrate-scale contact channel;
    \item $\tau_{\mathrm{D,bulk}}$ is the volumetric structural yield scale governing saturation of the complete active FCC cell.
\end{itemize}

The narrow contact footprint permits small-amplitude transverse motion while preserving a much larger volumetric failure threshold.

\subsection{Microscopic Surface-to-Volume Projection}

The calibrated contact cell defines a geometric transverse projection scale
\begin{equation}
    \iota_{\mathrm{tr}} \equiv \frac{A_{\mathrm{ct}}}{V_{\mathrm{p}}}.
\end{equation}

Substitution gives
\begin{equation}
    \boxed{ \iota_{\mathrm{tr}} = \frac{ \pi\ell_{\mathrm{S}}^2 }{ 2\sqrt{2}r_{\mathrm{D}}^3 } }.
\end{equation}

Its dimensions are
\begin{equation}
    [\iota_{\mathrm{tr}}] = \mathrm{m^{-1}}.
\end{equation}

The acoustic inertia density can therefore be expressed as
\begin{equation}
    \boxed{ \rho_{\mathrm{eff}}^{(0)} = \mu_{\mathrm{D}}\iota_{\mathrm{tr}} }.
\end{equation}

This is a microscopic surface-to-volume projection derived from FCC contact geometry and node-mass matching.

It must be distinguished from the localized-defect projection
\begin{equation}
    \rho_{\mathrm{IC}} = \mu_{\mathrm{D}}\iota_V
\end{equation}
used in the hydrodynamic mass law. The two quantities represent different reduced operations:
\begin{itemize}
    \item $\rho_{\mathrm{eff}}^{(0)}$ is the translational acoustic inertia density of the periodic primitive cell, whereas
    \item $\rho_{\mathrm{IC}}$ is the effective condensate impedance density sampled by a localized volumetric defect.
\end{itemize}

Equality of the two projections would require
\begin{equation}
    \iota_V = \iota_{\mathrm{tr}},
\end{equation}
which is not assumed by the present derivation. The contact-cell calculation therefore supplies the previously missing microscopic transverse projection without silently replacing the separately calibrated localized-defect density.

\subsection{Gravitationally Loaded Contact Cell}

The endpoint-transfer Jacobian in a locally uniform gravitational field is
\begin{equation}
    \mathcal{C} = \left( 1-\alpha_{\mathrm{G}} \right)^2.
\end{equation}

The completed contact-coordinate inertia is therefore
\begin{equation}
    m_{\mathrm{ct}}^{(\mathrm{eff})} = \frac{m_{\mathrm{ct}}}{\mathcal{C}^2}.
\end{equation}

Using
\begin{equation}
    m_{\mathrm{ct}} = M_{\mathrm{D}},
\end{equation}
one obtains
\begin{equation}
    \boxed{ m_{\mathrm{ct}}^{(\mathrm{eff})} = \frac{M_{\mathrm{D}}}{ \left( 1-\alpha_{\mathrm{G}} \right)^4 } }.
\end{equation}

The corresponding primitive-cell inertia density is
\begin{equation}
    \boxed{ \rho_{\mathrm{eff}} \left( \alpha_{\mathrm{G}} \right) = \frac{ \beta_{\mathrm{geo}}\rho_{\mathrm{D}} }{ \left( 1-\alpha_{\mathrm{G}} \right)^4 } }.
\end{equation}

The contact stiffness remains
\begin{equation}
    k_{\mathrm{ct}} = \frac{\pi}{2} \tau_{\mathrm{D}} \left( \frac{\ell_{\mathrm{S}}}{r_{\mathrm{D}}} \right)^2
\end{equation}
throughout the active FCC phase.

Therefore,
\begin{equation}
    c_P^2 = \frac{K_{\mathrm{sh}}}{\rho_{\mathrm{eff}}(\alpha_{\mathrm{G}})}.
\end{equation}

Using
\begin{equation}
    K_{\mathrm{sh}} = \beta_{\mathrm{geo}}\rho_{\mathrm{D}}c_{\mathrm{E}}^2,
\end{equation}
gives
\begin{equation}
    c_P^2 = c_{\mathrm{E}}^2 \left( 1-\alpha_{\mathrm{G}} \right)^4.
\end{equation}

Hence,
\begin{equation}
    \boxed{ c_P = c_{\mathrm{E}} \left( 1-\alpha_{\mathrm{G}} \right)^2 }.
\end{equation}

The gravitational propagation law is therefore recovered after absolute geometric and inertial calibration of the contact cell.

\subsection{Interstitial Volumetric Sector}

The contact-cell calibration fixes the transverse sector:
\begin{equation}
    A_{\mathrm{ct}}, \qquad m_{\mathrm{ct}}, \qquad k_{\mathrm{ct}}, \qquad \rho_{\mathrm{eff}}^{(0)}, \qquad K_{\mathrm{sh}}.
\end{equation}

It does not determine the independent tangent coefficient $B_{\mathrm{IC}}$ governing homogeneous compression of the interstitial condensate.

Let
\begin{equation}
    \epsilon_V = \operatorname{tr}\boldsymbol{\varepsilon}
\end{equation}
denote the primitive-cell volumetric strain, and let
\begin{equation}
    E_{\mathrm{IC}}(\epsilon_V)
\end{equation}
be the interstitial-condensate energy stored in one primitive cell.

The tangent condensate coefficient is defined by
\begin{equation}
    \boxed{ B_{\mathrm{IC}} = \frac{1}{V_{\mathrm{p}}} \left. \frac{\mathrm{d}^2E_{\mathrm{IC}}}{\mathrm{d}\epsilon_V^2} \right|_{\epsilon_V=0} }.
\end{equation}

Determination of $B_{\mathrm{IC}}$ requires a constitutive law for compression and redistribution of the octahedral and tetrahedral condensate regions. It cannot be obtained from the transverse surface-inertia matching alone.

The complete leading active-vacuum tensor remains
\begin{equation}
    C_{ijkl} = B_{\mathrm{IC}} \delta_{ij}\delta_{kl} + K_{\mathrm{sh}} \left( \delta_{ik}\delta_{jl} + \delta_{il}\delta_{jk} \right),
\end{equation}
with the transverse coefficient now absolutely calibrated.

\subsection{Constitutive Result}

The effective participating area of one nearest-neighbor FCC transfer coordinate is
\begin{equation}
    \boxed{ A_{\mathrm{ct}} = \frac{M_{\mathrm{D}}}{\mu_{\mathrm{D}}} = 2\pi\ell_{\mathrm{S}}^2 }.
\end{equation}

The corresponding fraction of one Wigner--Seitz face is
\begin{equation}
    \boxed{ f_{\mathrm{ct}} = \sqrt{2}\pi \left( \frac{\ell_{\mathrm{S}}}{r_{\mathrm{D}}} \right)^2 }.
\end{equation}

The integrated contact inertia is
\begin{equation}
    \boxed{ m_{\mathrm{ct}} = M_{\mathrm{D}} }.
\end{equation}

The integrated contact stiffness is
\begin{equation}
    \boxed{ k_{\mathrm{ct}} = \frac{\pi}{2} \tau_{\mathrm{D}} \left( \frac{\ell_{\mathrm{S}}}{r_{\mathrm{D}}} \right)^2 }.
\end{equation}

The unstrained primitive-cell acoustic inertia density is
\begin{equation}
    \boxed{ \rho_{\mathrm{eff}}^{(0)} = \frac{M_{\mathrm{D}}}{V_{\mathrm{p}}} = \beta_{\mathrm{geo}}\rho_{\mathrm{D}} }.
\end{equation}

The absolutely calibrated shear modulus is
\begin{equation}
    \boxed{ K_{\mathrm{sh}} = \frac{ A_{\mathrm{ct}}\tau_{\mathrm{D}} }{ V_{\mathrm{p}} } = \beta_{\mathrm{geo}} \rho_{\mathrm{D}} c_{\mathrm{E}}^2 }.
\end{equation}

These quantities satisfy
\begin{equation}
    \boxed{ \frac{K_{\mathrm{sh}}}{\rho_{\mathrm{eff}}^{(0)}} = \frac{\tau_{\mathrm{D}}}{\mu_{\mathrm{D}}} = c_{\mathrm{E}}^2 }.
\end{equation}

Under gravitational loading,
\begin{equation}
    \boxed{ \rho_{\mathrm{eff}} \left( \alpha_{\mathrm{G}} \right) = \frac{ \beta_{\mathrm{geo}}\rho_{\mathrm{D}} }{ \left( 1-\alpha_{\mathrm{G}} \right)^4 } },
\end{equation}
while
\begin{equation}
    \boxed{ K_{\mathrm{sh}} \left( \alpha_{\mathrm{G}} \right) = K_{\mathrm{sh}}(0) }.
\end{equation}

The resulting propagation law is
\begin{equation}
    \boxed{ c_P = c_{\mathrm{E}} \left( 1-\alpha_{\mathrm{G}} \right)^2 }.
\end{equation}

The transverse active-vacuum contact cell is therefore calibrated without an independent interface area, spring constant, nodal acoustic mass, or continuum shear normalization. The only unresolved leading linear coefficient is the interstitial-condensate tangent bulk response $B_{\mathrm{IC}}$.

\section{Interstitial-Condensate Constitutive Action and Longitudinal Response}

\subsection{Scope}

The preceding sections established the objective FCC contact action, the endpoint-transfer Jacobian, the twelve-neighbor homogenized shear tensor, and the absolute transverse stiffness--inertia normalization of the active vacuum. The remaining leading-order constitutive question concerns the response of the interstitial condensate to a local volumetric perturbation.

The active vacuum is a saturated dual-phase medium. Its primitive cell is partitioned into the solid FCC nodal fraction
\begin{equation}
    \beta_{\mathrm{geo}} = \frac{\pi}{3\sqrt{2}},
\end{equation}
and the interstitial-condensate fraction
\begin{equation}
    \beta_{\mathrm{IC}} = 1 - \beta_{\mathrm{geo}}.
\end{equation}

Hence,
\begin{equation}
    \boxed{ \beta_{\mathrm{geo}} + \beta_{\mathrm{IC}} = 1 }.
\end{equation}

No unoccupied third phase exists within the active FCC cell. The interstitial condensate continuously fills the connected octahedral and tetrahedral regions between the solid nodes.

The present analysis concerns small coherent perturbations about a quasistatically loaded active-vacuum state satisfying
\begin{equation}
    0 \leq \alpha_{\mathrm{G}} < 1,
\end{equation}
and
\begin{equation}
    \dot{\alpha}_{\mathrm{G}} \approx 0.
\end{equation}

The finite gravitational loading represented by $\alpha_{\mathrm{G}}$ is treated as a prescribed background state. The displacement field introduced below describes an incremental coherent disturbance about that state.

\subsection{Coherent Active-Vacuum Displacement}

Let
\begin{equation}
    \mathbf{u}(\mathbf{x},t)
\end{equation}
denote the homogenized displacement of the coherently coupled node--condensate system.

At the coherent-wave level, the FCC scaffold and the condensate do not constitute two independently propagating media. The condensate is entrained by the node-interface motion, and the already calibrated coefficient
\begin{equation}
    \rho_{\mathrm{eff}}^{(0)}
\end{equation}
represents the total effective inertia of that coupled transfer coordinate.

Accordingly, the coherent kinetic energy density is
\begin{equation}
    \boxed{ \mathcal{T} = \frac{1}{2} \rho_{\mathrm{eff}} \left( \alpha_{\mathrm{G}} \right) \dot{\mathbf{u}} \cdot \dot{\mathbf{u}} },
\end{equation}
where
\begin{equation}
    \boxed{ \rho_{\mathrm{eff}} \left( \alpha_{\mathrm{G}} \right) = \frac{ \rho_{\mathrm{eff}}^{(0)} }{ \left( 1-\alpha_{\mathrm{G}} \right)^4 } }.
\end{equation}

The condensate density used in the localized hydrodynamic mass law is not added independently to this coefficient. Relative condensate flow around a localized topological defect and coherent condensate entrainment in a propagating vacuum mode are distinct reductions of the same underlying medium. Adding both inertias to the same coherent coordinate would count the condensate response twice.

\subsection{Incremental Strain Decomposition}

The infinitesimal strain tensor is
\begin{equation}
    \varepsilon_{ij} = \frac{1}{2} \left( \partial_i u_j + \partial_j u_i \right).
\end{equation}

Its volumetric component is
\begin{equation}
    \boxed{ \varepsilon_V = \operatorname{tr}\boldsymbol{\varepsilon} = \nabla\cdot\mathbf{u} }.
\end{equation}

The deviatoric strain is
\begin{equation}
    \varepsilon_{ij}^{\mathrm{dev}} = \varepsilon_{ij} - \frac{1}{3} \varepsilon_V\delta_{ij}.
\end{equation}

The FCC contact-cell calibration supplies the shear coefficient
\begin{equation}
    K_{\mathrm{sh}} = \rho_{\mathrm{eff}}^{(0)} c_{\mathrm{E}}^2.
\end{equation}

Before imposing the condensate constraint, a general isotropic quadratic energy would therefore take the form
\begin{equation}
    \boxed{ \mathcal{W}_B = K_{\mathrm{sh}} \varepsilon_{ij}\varepsilon_{ij} + \frac{1}{2} B_{\mathrm{IC}} \varepsilon_V^2 },
\end{equation}
where $B_{\mathrm{IC}}$ is the coefficient multiplying the isotropic volumetric strain.

The associated elastic tensor is
\begin{equation}
    C_{ijkl} = B_{\mathrm{IC}} \delta_{ij}\delta_{kl} + K_{\mathrm{sh}} \left( \delta_{ik}\delta_{jl} + \delta_{il}\delta_{jk} \right).
\end{equation}

The corresponding total bulk modulus would be
\begin{equation}
    K_{\mathrm{bulk}} = B_{\mathrm{IC}} + \frac{2}{3}K_{\mathrm{sh}}.
\end{equation}

\subsection{Finite-Compressibility Regularization}

Consider first a finite value of $B_{\mathrm{IC}}$. The quadratic action is
\begin{equation}
    \boxed{ S_B^{(2)} = \int \left[ \frac{1}{2} \rho_{\mathrm{eff}} \dot{u}_i\dot{u}_i - K_{\mathrm{sh}} \varepsilon_{ij}\varepsilon_{ij} - \frac{1}{2} B_{\mathrm{IC}} \varepsilon_V^2 \right] \mathrm{d}^3x \, \mathrm{d}t }.
\end{equation}

The corresponding incremental stress is
\begin{equation}
    \boxed{ \sigma_{ij} = 2K_{\mathrm{sh}}\varepsilon_{ij} + B_{\mathrm{IC}} \varepsilon_V\delta_{ij} }.
\end{equation}

The field equation is
\begin{equation}
    \rho_{\mathrm{eff}}\ddot{u}_i = \partial_j\sigma_{ij},
\end{equation}
or
\begin{equation}
    \boxed{ \rho_{\mathrm{eff}}\ddot{\mathbf{u}} = K_{\mathrm{sh}}\nabla^2\mathbf{u} + \left( B_{\mathrm{IC}} + K_{\mathrm{sh}} \right) \nabla \left( \nabla\cdot\mathbf{u} \right) }.
\end{equation}

For a plane-wave perturbation,
\begin{equation}
    \mathbf{u} = \mathbf{e} \exp \left[ i \left( \mathbf{k}\cdot\mathbf{x}-\omega t \right) \right],
\end{equation}
the transverse branch satisfies
\begin{equation}
    \mathbf{k}\cdot\mathbf{e}_T = 0,
\end{equation}
and
\begin{equation}
    \boxed{ \omega_T^2 = \frac{ K_{\mathrm{sh}} }{ \rho_{\mathrm{eff}} } k^2 }.
\end{equation}

The longitudinal branch satisfies
\begin{equation}
    \mathbf{e}_L \parallel \mathbf{k},
\end{equation}
and
\begin{equation}
    \boxed{ \omega_L^2 = \frac{ B_{\mathrm{IC}} + 2K_{\mathrm{sh}} }{ \rho_{\mathrm{eff}} } k^2 }.
\end{equation}

Thus,
\begin{equation}
    c_T^2 = \frac{ K_{\mathrm{sh}} }{ \rho_{\mathrm{eff}} },
\end{equation}
while
\begin{equation}
    c_L^2 = \frac{ B_{\mathrm{IC}} + 2K_{\mathrm{sh}} }{ \rho_{\mathrm{eff}} }.
\end{equation}

\subsection{Exclusion of a Finite Longitudinal Acoustic Branch}

In the unstrained active vacuum,
\begin{equation}
    \rho_{\mathrm{eff}} = \rho_{\mathrm{eff}}^{(0)},
\end{equation}
and
\begin{equation}
    K_{\mathrm{sh}} = \rho_{\mathrm{eff}}^{(0)}c_{\mathrm{E}}^2.
\end{equation}

Therefore,
\begin{equation}
    c_T = c_{\mathrm{E}}.
\end{equation}

The longitudinal speed would be
\begin{equation}
    c_L^2 = c_{\mathrm{E}}^2 \left( \frac{ B_{\mathrm{IC}} }{ K_{\mathrm{sh}} } + 2 \right).
\end{equation}

Mechanical stability requires a positive total bulk modulus:
\begin{equation}
    B_{\mathrm{IC}} + \frac{2}{3}K_{\mathrm{sh}} > 0.
\end{equation}

Consequently,
\begin{equation}
    \frac{ B_{\mathrm{IC}} }{ K_{\mathrm{sh}} } > -\frac{2}{3}.
\end{equation}

It follows that every mechanically stable finite-compressibility branch satisfies
\begin{equation}
    c_L^2 > \frac{4}{3}c_{\mathrm{E}}^2,
\end{equation}
and therefore
\begin{equation}
    \boxed{ c_L > \sqrt{\frac{4}{3}} c_{\mathrm{E}} }.
\end{equation}

For a non-negative condensate compression coefficient,
\begin{equation}
    B_{\mathrm{IC}} \geq 0,
\end{equation}
the stronger bound is
\begin{equation}
    c_L \geq \sqrt{2}c_{\mathrm{E}}.
\end{equation}

A finite, stable, ordinary longitudinal acoustic branch would therefore exceed the established kinematic ceiling $c_{\mathrm{E}}$ of the active Soft Floor.

The coherent active-vacuum spectrum cannot contain such a branch.

The longitudinal sector must instead enter as a constrained response that supports static pressure and background loading without furnishing an independent finite-energy propagating polarization.

\subsection{Incremental Incompressibility of the Condensate-Filled Cell}

The required constitutive limit is
\begin{equation}
    \boxed{ B_{\mathrm{IC}} \rightarrow \infty }.
\end{equation}

For any finite-energy perturbation,
\begin{equation}
    \frac{1}{2} B_{\mathrm{IC}} \varepsilon_V^2 < \infty.
\end{equation}

Therefore,
\begin{equation}
    \varepsilon_V \rightarrow 0
\end{equation}
as $B_{\mathrm{IC}} \rightarrow \infty$.

The coherent active-vacuum displacement consequently satisfies
\begin{equation}
    \boxed{ \nabla\cdot\mathbf{u} = 0 }.
\end{equation}

This is incremental incompressibility about the prescribed background state. It does not assert that the active FCC phase cannot undergo quasistatic gravitational loading or eventual phase conversion. It states that a small coherent perturbation cannot create a propagating local volume fluctuation while the background value of $\alpha_{\mathrm{G}}$ remains fixed.

The distinction is
\begin{equation}
    \boxed{ \text{finite background compaction} \neq \text{incremental scalar wave} }.
\end{equation}

The background is parametrized by $\alpha_{\mathrm{G}}$. The coherent perturbation is volume preserving.

\subsection{Condensate Continuity}

Let
\begin{equation}
    \rho_{\mathrm{IC}}^{(0)}
\end{equation}
denote the equilibrium intrinsic density of the condensate within the interstitial network, and let
\begin{equation}
    \mathbf{v}_{\mathrm{IC}}
\end{equation}
be its coherent velocity.

Condensate mass conservation gives
\begin{equation}
    \frac{ \mathrm{D}\rho_{\mathrm{IC}} }{ \mathrm{D}t } + \rho_{\mathrm{IC}} \nabla\cdot\mathbf{v}_{\mathrm{IC}} = 0.
\end{equation}

Within the coherently entrained branch,
\begin{equation}
    \mathbf{v}_{\mathrm{IC}} = \dot{\mathbf{u}}.
\end{equation}

Incremental incompressibility of the condensate gives
\begin{equation}
    \frac{ \mathrm{D}\rho_{\mathrm{IC}} }{ \mathrm{D}t } = 0.
\end{equation}

Hence,
\begin{equation}
    \nabla\cdot\dot{\mathbf{u}} = 0.
\end{equation}

Integration about the unperturbed reference state gives
\begin{equation}
    \boxed{ \nabla\cdot\mathbf{u} = 0 }.
\end{equation}

The same constraint is therefore obtained from both the singular elastic limit and condensate mass conservation.

\subsection{Constrained Condensate Action}

Introduce the incremental condensate pressure
\begin{equation}
    \delta p_{\mathrm{IC}}(\mathbf{x},t)
\end{equation}
as a Lagrange multiplier enforcing the volume constraint.

The reduced coherent action is
\begin{equation}
    \boxed{ S_{\mathrm{IC}}^{(2)} = \int \left[ \frac{1}{2} \rho_{\mathrm{eff}} \left( \alpha_{\mathrm{G}} \right) \dot{u}_i\dot{u}_i - K_{\mathrm{sh}} \varepsilon_{ij}\varepsilon_{ij} + \delta p_{\mathrm{IC}} \nabla\cdot\mathbf{u} \right] \mathrm{d}^3x \, \mathrm{d}t }.
\end{equation}

Variation with respect to $\delta p_{\mathrm{IC}}$ gives
\begin{equation}
    \boxed{ \nabla\cdot\mathbf{u} = 0 }.
\end{equation}

Variation with respect to $\mathbf{u}$ gives
\begin{equation}
    \rho_{\mathrm{eff}} \left( \alpha_{\mathrm{G}} \right) \ddot{u}_i = \partial_j \delta\sigma_{ij},
\end{equation}
where the incremental stress is
\begin{equation}
    \boxed{ \delta\sigma_{ij} = - \delta p_{\mathrm{IC}} \delta_{ij} + 2K_{\mathrm{sh}} \varepsilon_{ij} }.
\end{equation}

Using the incompressibility condition,
\begin{equation}
    \partial_j\varepsilon_{ij} = \frac{1}{2}\nabla^2 u_i,
\end{equation}
the equation of motion becomes
\begin{equation}
    \boxed{ \rho_{\mathrm{eff}} \left( \alpha_{\mathrm{G}} \right) \ddot{\mathbf{u}} = K_{\mathrm{sh}} \nabla^2\mathbf{u} - \nabla\delta p_{\mathrm{IC}} }.
\end{equation}

The pressure field enforces the admissible displacement subspace but possesses no independent kinetic term.

\subsection{Pressure Constraint Equation}

Taking the divergence of the equation of motion gives
\begin{equation}
    \rho_{\mathrm{eff}} \partial_t^2 \left( \nabla\cdot\mathbf{u} \right) = K_{\mathrm{sh}} \nabla^2 \left( \nabla\cdot\mathbf{u} \right) - \nabla^2\delta p_{\mathrm{IC}}.
\end{equation}

Because
\begin{equation}
    \nabla\cdot\mathbf{u} = 0,
\end{equation}
one obtains
\begin{equation}
    \boxed{ \nabla^2\delta p_{\mathrm{IC}} = 0 }
\end{equation}
within a source-free homogeneous region.

In the presence of an applied incremental force density $\mathbf{f}$,
\begin{equation}
    \rho_{\mathrm{eff}}\ddot{\mathbf{u}} = K_{\mathrm{sh}}\nabla^2\mathbf{u} - \nabla\delta p_{\mathrm{IC}} + \mathbf{f},
\end{equation}
and the pressure satisfies
\begin{equation}
    \boxed{ \nabla^2\delta p_{\mathrm{IC}} = \nabla\cdot\mathbf{f} }.
\end{equation}

The pressure equation is a constraint equation determined by the applied loading and boundary conditions. It is not a hyperbolic wave equation and does not define an independent propagating scalar degree of freedom.

\subsection{Absence of a Longitudinal Eigenmode}

For a plane-wave perturbation, the constraint requires
\begin{equation}
    \mathbf{k}\cdot\mathbf{e} = 0.
\end{equation}

A longitudinal polarization would satisfy
\begin{equation}
    \mathbf{e}_L \parallel \mathbf{k},
\end{equation}
and therefore
\begin{equation}
    \mathbf{k}\cdot\mathbf{e}_L \neq 0.
\end{equation}

Such a polarization is excluded from the admissible finite-energy configuration space:
\begin{equation}
    \boxed{ \mathbf{e}_L \notin \mathcal{H}_{\mathrm{active}} },
\end{equation}
where $\mathcal{H}_{\mathrm{active}}$ denotes the linearized coherent-mode space of the active FCC phase.

The formal longitudinal frequency of the finite-$B_{\mathrm{IC}}$ regularization obeys
\begin{equation}
    \omega_L^2 = \frac{ B_{\mathrm{IC}} + 2K_{\mathrm{sh}} }{ \rho_{\mathrm{eff}} } k^2.
\end{equation}

As $B_{\mathrm{IC}} \rightarrow \infty$, one has
\begin{equation}
    \omega_L \rightarrow \infty
\end{equation}
for every fixed nonzero $k$. The longitudinal excitation is therefore removed from the finite-frequency spectrum rather than converted into a superluminal propagating signal.

The remaining coherent acoustic spectrum contains only the two transverse polarizations.

\subsection{Transverse Dispersion}

For an admissible transverse plane wave,
\begin{equation}
    \mathbf{k}\cdot\mathbf{e}_T = 0.
\end{equation}

Taking the scalar product of the field equation with $\mathbf{e}_T$ removes the pressure term because
\begin{equation}
    \mathbf{e}_T\cdot\mathbf{k} = 0.
\end{equation}

The dispersion relation is
\begin{equation}
    \boxed{ \omega_T^2 = \frac{ K_{\mathrm{sh}} }{ \rho_{\mathrm{eff}} \left( \alpha_{\mathrm{G}} \right) } k^2 + \mathcal{O}\left( k^4d_{\mathrm{nn}}^2 \right) }.
\end{equation}

Using
\begin{equation}
    K_{\mathrm{sh}} = \rho_{\mathrm{eff}}^{(0)} c_{\mathrm{E}}^2,
\end{equation}
and
\begin{equation}
    \rho_{\mathrm{eff}} \left( \alpha_{\mathrm{G}} \right) = \frac{ \rho_{\mathrm{eff}}^{(0)} }{ \left( 1-\alpha_{\mathrm{G}} \right)^4 },
\end{equation}
gives
\begin{equation}
    \omega_T^2 = c_{\mathrm{E}}^2 \left( 1-\alpha_{\mathrm{G}} \right)^4 k^2 + \mathcal{O}\left( k^4d_{\mathrm{nn}}^2 \right).
\end{equation}

Therefore,
\begin{equation}
    \boxed{ c_P = c_{\mathrm{E}} \left( 1-\alpha_{\mathrm{G}} \right)^2 }.
\end{equation}

Both transverse polarizations retain the same speed because the leading homogenized shear tensor is isotropic.

\subsection{Static Longitudinal Pressure}

The absence of a propagating longitudinal mode does not imply the absence of longitudinal stress.

The total condensate pressure may be decomposed as
\begin{equation}
    p_{\mathrm{IC}} = \overline{p}_{\mathrm{IC}} + \delta p_{\mathrm{IC}},
\end{equation}
where:
\begin{itemize}
    \item $\overline{p}_{\mathrm{IC}}$ is the quasistatic background pressure associated with the prescribed gravitational loading;
    \item $\delta p_{\mathrm{IC}}$ is the incremental reaction pressure enforcing coherent incompressibility.
\end{itemize}

The structural gravitational loading is
\begin{equation}
    \tau_L = \alpha_{\mathrm{G}}\tau_{\mathrm{D}}.
\end{equation}

Its volumetric pressure equivalent is obtained by projection through the node radius:
\begin{equation}
    \overline{p}_{\mathrm{IC}} = \frac{\tau_L}{r_{\mathrm{D}}}.
\end{equation}

Using
\begin{equation}
    \tau_{\mathrm{D,bulk}} = \frac{\tau_{\mathrm{D}}}{r_{\mathrm{D}}},
\end{equation}
one obtains
\begin{equation}
    \boxed{ \overline{p}_{\mathrm{IC}} = \alpha_{\mathrm{G}} \tau_{\mathrm{D,bulk}} }.
\end{equation}

Thus, the background longitudinal pressure increases continuously with the occupied gravitational capacity.

At zero gravitational loading,
\begin{equation}
    \alpha_{\mathrm{G}} = 0,
\end{equation}
and
\begin{equation}
    \overline{p}_{\mathrm{IC}} = 0.
\end{equation}

At the active-phase boundary,
\begin{equation}
    \alpha_{\mathrm{G}} \rightarrow 1^{-},
\end{equation}
the pressure approaches
\begin{equation}
    \boxed{ \overline{p}_{\mathrm{IC}} \rightarrow \tau_{\mathrm{D,bulk}} }.
\end{equation}

The finite quantity $\tau_{\mathrm{D,bulk}}$ is therefore the saturation pressure at which the active FCC phase terminates. It is not the finite tangent modulus of an ordinary compressible acoustic branch.

\subsection{Distinction Between Constraint Stiffness and Yield Pressure}

The constitutive roles of $B_{\mathrm{IC}}$ and $\tau_{\mathrm{D,bulk}}$ are distinct.

The coefficient $B_{\mathrm{IC}}$ appears in the finite-compressibility regularization:
\begin{equation}
    \frac{1}{2} B_{\mathrm{IC}} \varepsilon_V^2.
\end{equation}

The coherent active-vacuum limit requires
\begin{equation}
    \boxed{ B_{\mathrm{IC}} \rightarrow \infty },
\end{equation}
so that
\begin{equation}
    \varepsilon_V = 0.
\end{equation}

By contrast, $\tau_{\mathrm{D,bulk}}$ is the finite pressure threshold of the active FCC phase:
\begin{equation}
    \boxed{ \overline{p}_{\mathrm{IC}} \leq \tau_{\mathrm{D,bulk}}, \qquad 0 \leq \alpha_{\mathrm{G}} < 1 }.
\end{equation}

At equality,
\begin{equation}
    \overline{p}_{\mathrm{IC}} = \tau_{\mathrm{D,bulk}},
\end{equation}
the FCC coordinates cease to describe the local state and the structure transitions to the phase-locked SC substrate.

The limit $B_{\mathrm{IC}} \rightarrow \infty$ therefore does not imply an infinite realized pressure. It means that finite coherent volume strain is excluded. The realized pressure remains a reaction field bounded by the phase-transition threshold.

This is the same mechanical distinction found in an incrementally incompressible medium: the bulk constraint can be singular while the hydrostatic pressure remains finite.

\subsection{Gravitationally Preloaded Incremental Dynamics}

The background pressure modifies the local transfer capacity through
\begin{equation}
    \alpha_{\mathrm{G}} = \frac{ \overline{p}_{\mathrm{IC}} }{ \tau_{\mathrm{D,bulk}} }.
\end{equation}

The endpoint-transfer Jacobian is
\begin{equation}
    \mathcal{C} = \left( 1-\alpha_{\mathrm{G}} \right)^2.
\end{equation}

Consequently, the background longitudinal pressure changes the incremental transverse inertia:
\begin{equation}
    \rho_{\mathrm{eff}} = \frac{ \rho_{\mathrm{eff}}^{(0)} }{ \mathcal{C}^2 }.
\end{equation}

The background pressure does not generate a new scalar wave. It modifies the transverse spectrum by reducing the completed-transfer Jacobian.

The mechanical sequence is therefore
\begin{equation}
    \boxed{ \overline{p}_{\mathrm{IC}} \longrightarrow \alpha_{\mathrm{G}} \longrightarrow \mathcal{C} \longrightarrow \rho_{\mathrm{eff}} \longrightarrow c_P }.
\end{equation}

Explicitly,
\begin{equation}
    \overline{p}_{\mathrm{IC}} = \alpha_{\mathrm{G}} \tau_{\mathrm{D,bulk}},
\end{equation}
\begin{equation}
    \mathcal{C} = \left( 1-\alpha_{\mathrm{G}} \right)^2,
\end{equation}
\begin{equation}
    \rho_{\mathrm{eff}} = \frac{ \rho_{\mathrm{eff}}^{(0)} }{ \left( 1-\alpha_{\mathrm{G}} \right)^4 },
\end{equation}
and
\begin{equation}
    c_P = c_{\mathrm{E}} \left( 1-\alpha_{\mathrm{G}} \right)^2.
\end{equation}

The static longitudinal sector and the dynamic transverse sector are thus mechanically coupled without becoming the same mode.

\subsection{Relative Condensate Motion}

The constrained coherent action applies to freely propagating perturbations of the active vacuum. It does not eliminate relative condensate motion in the presence of a localized topological defect.

For a defect moving through the lattice, the condensate velocity need not equal the coherent scaffold velocity. The relative flow produces the hydrodynamic impedance represented by
\begin{equation}
    m = \rho_{\mathrm{IC}} V_{\mathrm{eff}} \sin(\Delta\varphi).
\end{equation}

That response is localized, defect-driven, and generally nonuniform. It belongs to the relative-flow sector of the active vacuum rather than to the homogeneous acoustic spectrum derived here.

The distinction is
\begin{equation}
    \boxed{
    \begin{aligned}
        \text{coherent vacuum mode:}&\quad \mathbf{v}_{\mathrm{IC}} = \dot{\mathbf{u}}, \\[3pt]
        \text{localized defect response:}&\quad \mathbf{v}_{\mathrm{IC}} - \dot{\mathbf{u}} \neq \mathbf{0}.
    \end{aligned}
    }
\end{equation}

The first produces transverse propagation. The second produces inertial impedance, damping, and inductive condensate response.

A separate relative-flow action may be introduced when the localized defect sector is derived. It is not an additional free scalar acoustic branch of the unperturbed active vacuum.

\subsection{Constitutive Result}

The interstitial-condensate sector of the coherent active vacuum is incrementally incompressible:
\begin{equation}
    \boxed{ \nabla\cdot\mathbf{u} = 0 }.
\end{equation}

The finite-compressibility coefficient introduced in the preceding homogenization is a regularization parameter whose active-vacuum limit is
\begin{equation}
    \boxed{ B_{\mathrm{IC}} \rightarrow \infty }.
\end{equation}

The coherent quadratic action is
\begin{equation}
    \boxed{ S_{\mathrm{IC}}^{(2)} = \int \left[ \frac{1}{2} \rho_{\mathrm{eff}} \left( \alpha_{\mathrm{G}} \right) \dot{u}_i\dot{u}_i - K_{\mathrm{sh}} \varepsilon_{ij}\varepsilon_{ij} + \delta p_{\mathrm{IC}} \nabla\cdot\mathbf{u} \right] \mathrm{d}^3x \, \mathrm{d}t }.
\end{equation}

Its incremental stress is
\begin{equation}
    \boxed{ \delta\sigma_{ij} = - \delta p_{\mathrm{IC}}\delta_{ij} + 2K_{\mathrm{sh}}\varepsilon_{ij} }.
\end{equation}

The field equations are
\begin{equation}
    \boxed{ \rho_{\mathrm{eff}} \left( \alpha_{\mathrm{G}} \right) \ddot{\mathbf{u}} = K_{\mathrm{sh}} \nabla^2\mathbf{u} - \nabla\delta p_{\mathrm{IC}} },
\end{equation}
and
\begin{equation}
    \boxed{ \nabla\cdot\mathbf{u} = 0 }.
\end{equation}

No finite-frequency longitudinal acoustic eigenmode exists in the coherent active-vacuum spectrum.

The static background pressure is
\begin{equation}
    \boxed{ \overline{p}_{\mathrm{IC}} = \alpha_{\mathrm{G}} \tau_{\mathrm{D,bulk}} },
\end{equation}
and reaches the FCC saturation threshold at
\begin{equation}
    \boxed{ \overline{p}_{\mathrm{IC}} = \tau_{\mathrm{D,bulk}} \qquad \text{when} \qquad \alpha_{\mathrm{G}} = 1 }.
\end{equation}

The two admissible coherent acoustic branches are transverse and satisfy
\begin{equation}
    \boxed{ \omega_T^2 = c_{\mathrm{E}}^2 \left( 1-\alpha_{\mathrm{G}} \right)^4 k^2 + \mathcal{O}\left( k^4d_{\mathrm{nn}}^2 \right) },
\end{equation}
with propagation speed
\begin{equation}
    \boxed{ c_P = c_{\mathrm{E}} \left( 1-\alpha_{\mathrm{G}} \right)^2 }.
\end{equation}

The interstitial condensate therefore performs three distinct mechanical functions:
\begin{equation}
    \boxed{ \text{incremental volume constraint} },
\end{equation}
\begin{equation}
    \boxed{ \text{static longitudinal pressure support} },
\end{equation}
and
\begin{equation}
    \boxed{ \text{localized hydrodynamic impedance under relative defect motion} }.
\end{equation}

It does not introduce an independent propagating scalar mode into the coherent active-vacuum spectrum.

\section{Localized Defect--Condensate Relative-Flow Action and Microscopic Impedance Projection}

\subsection{Scope}

The preceding section established the coherent active-vacuum sector in which the FCC scaffold and the interstitial condensate move as an entrained, incrementally incompressible system. A localized topological defect introduces a distinct response. During translation, internal reconfiguration, or phase transfer, the condensate surrounding the defect need not remain comoving with the local scaffold. The resulting relative flow stores kinetic energy and produces an added inertial impedance.

The present section formulates this localized relative-flow sector and gives a mechanical definition of the surface-to-volume projection previously represented by
\begin{equation}
    \rho_{\mathrm{IC}} = \mu_{\mathrm{D}}\iota_V.
\end{equation}

The projection is derived as a mode-weighted ratio between the defect interface and the condensate volume participating in the induced flow. It is therefore determined by the defect geometry, boundary motion, interstitial connectivity, and condensate response field rather than by an externally selected dimensional conversion.

The analysis is restricted to the conservative low-frequency branch. Dissipative boundary layers, irreversible heating, and radiative loss are excluded at the retained order.

\subsection{Relative Scaffold--Condensate Kinematics}

Let
\begin{equation}
    \mathbf{u}(\mathbf{x},t)
\end{equation}
denote the local FCC scaffold displacement, and let
\begin{equation}
    \mathbf{w}(\mathbf{x},t)
\end{equation}
denote the displacement of the interstitial condensate.

The condensate velocity relative to the scaffold is
\begin{equation}
    \boxed{ \mathbf{v}_{\mathrm{rel}} = \dot{\mathbf{w}} - \dot{\mathbf{u}} }.
\end{equation}

For a coherent transverse vacuum mode,
\begin{equation}
    \mathbf{v}_{\mathrm{rel}} = \mathbf{0}.
\end{equation}

For a localized defect,
\begin{equation}
    \mathbf{v}_{\mathrm{rel}} \neq \mathbf{0}
\end{equation}
within the surrounding connected interstitial domain.

Let $q(t)$ be a local generalized coordinate describing the selected defect motion. The variable $q$ may represent translation of the defect core, a collective boundary displacement, or another localized configuration coordinate. Its definition is confined to the present derivation.

For a linearized defect mode, write
\begin{equation}
    \boxed{ \mathbf{v}_{\mathrm{rel}}(\mathbf{x},t) = \dot{q}(t) \boldsymbol{\Upsilon}(\mathbf{x}) },
\end{equation}
where
\begin{equation}
    \boldsymbol{\Upsilon}(\mathbf{x})
\end{equation}
is the dimensionless condensate response field generated by unit generalized defect velocity.

\subsection{Defect Boundary Kinematics}

Let
\begin{equation}
    \partial\Omega_{\mathrm{def}}
\end{equation}
denote the effective interface between the localized defect and the surrounding interstitial-condensate domain
\begin{equation}
    \Omega_{\mathrm{IC}}^{(\mathrm{def})}.
\end{equation}

The position of a material point on the defect boundary is written as
\begin{equation}
    \mathbf{X}_{\partial} = \mathbf{X}_{\partial}(q).
\end{equation}

Define the dimensionless boundary mode by
\begin{equation}
    \boxed{ \mathbf{b}_{\partial} = \frac{\partial\mathbf{X}_{\partial}}{\partial q} }.
\end{equation}

The boundary velocity relative to the local scaffold is then
\begin{equation}
    \dot{\mathbf{X}}_{\partial}^{(\mathrm{rel})} = \dot{q} \mathbf{b}_{\partial}.
\end{equation}

At conservative order, the defect and condensate satisfy the no-penetration condition
\begin{equation}
    \boxed{ \boldsymbol{\Upsilon}\cdot\hat{\mathbf{n}} = \mathbf{b}_{\partial}\cdot\hat{\mathbf{n}} \qquad \text{on } \partial\Omega_{\mathrm{def}} },
\end{equation}
where $\hat{\mathbf{n}}$ is the unit normal directed from the defect into the condensate.

This boundary condition requires the normal condensate velocity to follow the defect boundary while permitting nondissipative tangential slip. A viscous no-slip condition, when required, belongs to the dissipative extension of the relative-flow theory.

\subsection{Conservative Relative-Flow Boundary-Value Problem}

The interstitial condensate is incrementally incompressible within the active FCC phase. The localized response field therefore satisfies
\begin{equation}
    \boxed{ \nabla\cdot\boldsymbol{\Upsilon} = 0 }
\end{equation}
throughout
\begin{equation}
    \Omega_{\mathrm{IC}}^{(\mathrm{def})}.
\end{equation}

On the zero-circulation conservative branch, the minimum-energy flow is also irrotational:
\begin{equation}
    \boxed{ \nabla\times\boldsymbol{\Upsilon} = \mathbf{0} }.
\end{equation}

A nonzero conserved circulation, if supported by a chiral or vortex-like topological sector, must be treated as an additional branch rather than included implicitly in the minimum-energy response.

For an isolated defect,
\begin{equation}
    \boldsymbol{\Upsilon} \rightarrow \mathbf{0}
\end{equation}
as the distance from the defect tends to infinity.

For a defect embedded in a periodic FCC calculation, $\boldsymbol{\Upsilon}$ satisfies periodic boundary conditions and zero net relative flux through the primitive-cell boundary.

The incompressibility compatibility condition is
\begin{equation}
    \boxed{ \int_{\partial\Omega_{\mathrm{IC}}^{(\mathrm{def})}} \boldsymbol{\Upsilon}\cdot\hat{\mathbf{n}} \, \mathrm{d}A = 0 }.
\end{equation}

Equivalently, the conservative response field is the minimizer of
\begin{equation}
    \int_{\Omega_{\mathrm{IC}}^{(\mathrm{def})}} \left| \boldsymbol{\Upsilon} \right|^2 \, \mathrm{d}V
\end{equation}
subject to the incompressibility condition and the prescribed normal boundary motion.

This variational problem defines the added-inertia response for a specified defect geometry and boundary mode.

\subsection{Normalization Invariance}

The normalization of the generalized coordinate $q$ is arbitrary. Under the rescaling
\begin{equation}
    q^{\prime} = aq,
\end{equation}
the boundary and condensate modes transform as
\begin{equation}
    \mathbf{b}_{\partial}^{\prime} = \frac{1}{a}\mathbf{b}_{\partial},
\end{equation}
and
\begin{equation}
    \boldsymbol{\Upsilon}^{\prime} = \frac{1}{a}\boldsymbol{\Upsilon}.
\end{equation}

The corresponding mode integrals both scale as $a^{-2}$, while
\begin{equation}
    \dot{q}^{\prime 2} = a^2\dot{q}^2.
\end{equation}

The kinetic energy and the ratio of the mode integrals are therefore invariant under the normalization of $q$. The microscopic impedance projection derived below is consequently a property of the physical defect mode rather than of the chosen generalized-coordinate scale.

\subsection{Mode-Weighted Interface Area}

Define the effective interface measure associated with the defect mode by
\begin{equation}
    \boxed{ A_{\mathrm{imp}} = \int_{\partial\Omega_{\mathrm{def}}} \left| \mathbf{b}_{\partial} \right|^2 \, \mathrm{d}A }.
\end{equation}

Because $\mathbf{b}_{\partial}$ is dimensionless,
\begin{equation}
    [A_{\mathrm{imp}}] = \mathrm{m}^2.
\end{equation}

The quantity $A_{\mathrm{imp}}$ is not generally equal to the total geometric area of the defect boundary. It is weighted by the displacement produced by the selected collective mode.

Boundary regions that remain stationary do not contribute. Regions undergoing larger displacement relative to $q$ receive proportionally greater weight.

\subsection{Mode-Weighted Condensate Volume}

Define the condensate participation volume by
\begin{equation}
    \boxed{ V_{\mathrm{imp}} = \int_{\Omega_{\mathrm{IC}}^{(\mathrm{def})}} \left| \boldsymbol{\Upsilon} \right|^2 \, \mathrm{d}V }.
\end{equation}

Its dimensions are
\begin{equation}
    [V_{\mathrm{imp}}] = \mathrm{m}^3.
\end{equation}

This is an inertia-weighted volume rather than a purely geometric volume. A condensate region with negligible relative velocity contributes negligibly even when it is geometrically adjacent to the defect. Conversely, an extended flow of smaller amplitude may contribute appreciably when integrated over a sufficiently large domain.

The kinetic participation volume is therefore determined by the solved response field and is not automatically identical to the geometric volume swept by the defect.

\subsection{Phase-Participation Function}

Let
\begin{equation}
    f_{\varphi}(\Delta\varphi)
\end{equation}
denote the fraction of the available interface inertia activated by the local phase mismatch.

Because $f_{\varphi}$ multiplies a quadratic kinetic term, it must satisfy
\begin{equation}
    \boxed{ f_{\varphi}(\Delta\varphi) \geq 0 }.
\end{equation}

For an unpolarized reciprocal vacuum, the participation function is even:
\begin{equation}
    \boxed{ f_{\varphi}(-\Delta\varphi) = f_{\varphi}(\Delta\varphi) }.
\end{equation}

Phase periodicity requires
\begin{equation}
    \boxed{ f_{\varphi} \left( \Delta\varphi+2\pi \right) = f_{\varphi}(\Delta\varphi) }.
\end{equation}

A phase-matched state that produces no localized hydrodynamic obstruction satisfies
\begin{equation}
    \boxed{ f_{\varphi}(0) = 0 }.
\end{equation}

The precise functional form of $f_{\varphi}$ must follow from the nonlinear phase-contact action. It is not fixed by the surface-to-volume projection.

The reduced treatment below assumes that the phase mismatch is uniform over the active defect mode. For a spatially varying phase mismatch, $f_{\varphi}$ must remain inside the corresponding surface and volume integrals.

\subsection{Surface Representation of the Localized Kinetic Energy}

The active-vacuum surface inertia is $\mu_{\mathrm{D}}$. The kinetic energy of the selected localized defect mode may therefore be represented in interface form as
\begin{equation}
    \boxed{ T_{\partial} = \frac{1}{2} \mu_{\mathrm{D}} A_{\mathrm{imp}} f_{\varphi}(\Delta\varphi) \dot{q}^2 }.
\end{equation}

The associated added-inertia coefficient is
\begin{equation}
    \boxed{ m_{\partial} = \mu_{\mathrm{D}} A_{\mathrm{imp}} f_{\varphi}(\Delta\varphi) }.
\end{equation}

This expression represents the inertial loading of the active boundary mode measured through the established surface coefficient $\mu_{\mathrm{D}}$.

\subsection{Volumetric Representation of the Same Mode}

The same kinetic energy may be expressed through an effective volumetric condensate impedance:
\begin{equation}
    \boxed{ T_{\mathrm{vol}} = \frac{1}{2} \rho_{\mathrm{IC}}^{(\mathrm{imp})} V_{\mathrm{imp}} f_{\varphi}(\Delta\varphi) \dot{q}^2 },
\end{equation}
where
\begin{equation}
    \rho_{\mathrm{IC}}^{(\mathrm{imp})}
\end{equation}
is the effective volumetric impedance density sampled by the selected defect mode.

The surface and volumetric representations describe the same kinetic energy. Therefore,
\begin{equation}
    T_{\partial} = T_{\mathrm{vol}}.
\end{equation}

For a nonzero phase-participation factor,
\begin{equation}
    \mu_{\mathrm{D}}A_{\mathrm{imp}} = \rho_{\mathrm{IC}}^{(\mathrm{imp})} V_{\mathrm{imp}}.
\end{equation}

Hence,
\begin{equation}
    \boxed{ \rho_{\mathrm{IC}}^{(\mathrm{imp})} = \mu_{\mathrm{D}} \frac{A_{\mathrm{imp}}}{V_{\mathrm{imp}}} }.
\end{equation}

\subsection{Microscopic Impedance Projection}

Define the defect-mode surface-to-volume projection by
\begin{equation}
    \boxed{ \iota_V^{(\mathrm{def})} \equiv \frac{A_{\mathrm{imp}}}{V_{\mathrm{imp}}} }.
\end{equation}

Its dimensions are
\begin{equation}
    \left[ \iota_V^{(\mathrm{def})} \right] = \mathrm{m}^{-1}.
\end{equation}

The effective condensate impedance density is consequently
\begin{equation}
    \boxed{ \rho_{\mathrm{IC}}^{(\mathrm{imp})} = \mu_{\mathrm{D}} \iota_V^{(\mathrm{def})} }.
\end{equation}

Using the defining mode integrals,
\begin{equation}
    \boxed{ \iota_V^{(\mathrm{def})} = \frac{ \displaystyle \int_{\partial\Omega_{\mathrm{def}}} \left| \mathbf{b}_{\partial} \right|^2 \, \mathrm{d}A }{ \displaystyle \int_{\Omega_{\mathrm{IC}}^{(\mathrm{def})}} \left| \boldsymbol{\Upsilon} \right|^2 \, \mathrm{d}V } }.
\end{equation}

The projection is therefore a mode-weighted inverse length determined by the localized relative-flow solution.

It is invariant under rescaling of the generalized coordinate and depends only on the physical defect geometry, boundary conditions, and condensate response.

\subsection{Impedance Penetration Length}

Define the associated impedance penetration length by
\begin{equation}
    \boxed{ \ell_{\mathrm{imp}}^{(\mathrm{def})} \equiv \frac{ V_{\mathrm{imp}} }{ A_{\mathrm{imp}} } }.
\end{equation}

Then,
\begin{equation}
    \boxed{ \iota_V^{(\mathrm{def})} = \frac{1}{ \ell_{\mathrm{imp}}^{(\mathrm{def})} } }.
\end{equation}

The projected density becomes
\begin{equation}
    \boxed{ \rho_{\mathrm{IC}}^{(\mathrm{imp})} = \frac{ \mu_{\mathrm{D}} }{ \ell_{\mathrm{imp}}^{(\mathrm{def})} } }.
\end{equation}

The quantity previously represented by an unspecified projection factor is therefore the inverse penetration length of the condensate flow induced by the localized defect.

The empirical assignment
\begin{equation}
    \iota_V = 1 \ \mathrm{m}^{-1}
\end{equation}
corresponds to
\begin{equation}
    \ell_{\mathrm{imp}} = 1 \ \mathrm{m}.
\end{equation}

Such a value can be retained only as an empirical calibration until it is produced by a solved defect-flow boundary-value problem.

\subsection{Added Inertia of the Localized Defect}

The kinetic energy of the generalized defect coordinate has the standard form
\begin{equation}
    T_{\mathrm{rel}} = \frac{1}{2} m_{\mathrm{add}} \dot{q}^2.
\end{equation}

Comparison with the surface representation gives
\begin{equation}
    \boxed{ m_{\mathrm{add}} = \mu_{\mathrm{D}} A_{\mathrm{imp}} f_{\varphi}(\Delta\varphi) }.
\end{equation}

Comparison with the volumetric representation gives
\begin{equation}
    \boxed{ m_{\mathrm{add}} = \rho_{\mathrm{IC}}^{(\mathrm{imp})} V_{\mathrm{imp}} f_{\varphi}(\Delta\varphi) }.
\end{equation}

The two expressions are identical because
\begin{equation}
    \rho_{\mathrm{IC}}^{(\mathrm{imp})} V_{\mathrm{imp}} = \mu_{\mathrm{D}} A_{\mathrm{imp}}.
\end{equation}

The inertial mass is therefore the quadratic kinetic coefficient generated by the defect-induced condensate response. It is not the static material mass of the defect core.

\subsection{Geometric Swept-Volume Reduction}

Consider a localized defect whose active boundary region occupies a thin spherical shell surrounding a node of radius $r_{\mathrm{D}}$. Let
\begin{equation}
    r_{\mathrm{amp}}
\end{equation}
denote the equilibrium radial thickness of the active defect structure.

The corresponding geometric shell volume is
\begin{equation}
    V_{\mathrm{sweep}} = \frac{4\pi}{3} \left[ \left( r_{\mathrm{D}}+r_{\mathrm{amp}} \right)^3 - r_{\mathrm{D}}^3 \right].
\end{equation}

Expansion gives
\begin{equation}
    V_{\mathrm{sweep}} = 4\pi r_{\mathrm{D}}^2 r_{\mathrm{amp}} + 4\pi r_{\mathrm{D}}r_{\mathrm{amp}}^2 + \frac{4\pi}{3}r_{\mathrm{amp}}^3.
\end{equation}

For
\begin{equation}
    r_{\mathrm{amp}} \ll r_{\mathrm{D}},
\end{equation}
the leading contribution is
\begin{equation}
    \boxed{ V_{\mathrm{sweep}} = 4\pi r_{\mathrm{D}}^2 r_{\mathrm{amp}} + \mathcal{O} \left( r_{\mathrm{D}}r_{\mathrm{amp}}^2 \right) }.
\end{equation}

This gives the geometric origin of the reduced expression
\begin{equation}
    V_{\mathrm{eff}} = 4\pi r_{\mathrm{D}}^2 r_{\mathrm{amp}}.
\end{equation}

The inertially participating volume remains
\begin{equation}
    V_{\mathrm{imp}} = \int_{\Omega_{\mathrm{IC}}^{(\mathrm{def})}} \left| \boldsymbol{\Upsilon} \right|^2 \, \mathrm{d}V,
\end{equation}
and need not equal $V_{\mathrm{sweep}}$.

Define the dimensionless added-inertia geometry factor by
\begin{equation}
    \boxed{ \mathcal{G}_{\mathrm{imp}} \equiv \frac{ V_{\mathrm{imp}} }{ V_{\mathrm{sweep}} } }.
\end{equation}

Therefore,
\begin{equation}
    \boxed{ V_{\mathrm{imp}} = \mathcal{G}_{\mathrm{imp}} V_{\mathrm{sweep}} }.
\end{equation}

In the thin-shell limit,
\begin{equation}
    \boxed{ V_{\mathrm{imp}} = 4\pi \mathcal{G}_{\mathrm{imp}} r_{\mathrm{D}}^2 r_{\mathrm{amp}} + \mathcal{O} \left( r_{\mathrm{D}}r_{\mathrm{amp}}^2 \right) }.
\end{equation}

The condition
\begin{equation}
    \mathcal{G}_{\mathrm{imp}} = 1
\end{equation}
is a property of the induced condensate flow. It is not implied by the geometric shell construction alone.

\subsection{Reduced Hydrodynamic Mass Law}

Substitution of the thin-shell participation volume into the added-inertia expression gives
\begin{equation}
    m_{\mathrm{add}} = 4\pi r_{\mathrm{D}}^2 \rho_{\mathrm{IC}}^{(\mathrm{imp})} \mathcal{G}_{\mathrm{imp}} r_{\mathrm{amp}} f_{\varphi}(\Delta\varphi) + \mathcal{O} \left( r_{\mathrm{D}}r_{\mathrm{amp}}^2 \right).
\end{equation}

Define the mode-resolved impedance coefficient by
\begin{equation}
    \boxed{ K_{\mathrm{D}}^{(\mathrm{def})} \equiv 4\pi r_{\mathrm{D}}^2 \rho_{\mathrm{IC}}^{(\mathrm{imp})} \mathcal{G}_{\mathrm{imp}} }.
\end{equation}

Its dimensions are
\begin{equation}
    \left[ K_{\mathrm{D}}^{(\mathrm{def})} \right] = \mathrm{kg\,m^{-1}}.
\end{equation}

The leading reduced mass law is therefore
\begin{equation}
    \boxed{ m_{\mathrm{add}} = K_{\mathrm{D}}^{(\mathrm{def})} r_{\mathrm{amp}} f_{\varphi}(\Delta\varphi) }.
\end{equation}

Using the microscopic projection,
\begin{equation}
    \boxed{ K_{\mathrm{D}}^{(\mathrm{def})} = 4\pi r_{\mathrm{D}}^2 \mu_{\mathrm{D}} \iota_V^{(\mathrm{def})} \mathcal{G}_{\mathrm{imp}} }.
\end{equation}

The previously used impedance constant
\begin{equation}
    K_{\mathrm{D}} = 4\pi r_{\mathrm{D}}^2\rho_{\mathrm{IC}}
\end{equation}
is recovered when the projected density is universal within the relevant defect class and
\begin{equation}
    \mathcal{G}_{\mathrm{imp}} = 1.
\end{equation}

\subsection{Status of the Sine Phase Closure}

The earlier reduced mass law contains the phase factor
\begin{equation}
    \sin(\Delta\varphi).
\end{equation}

The present relative-flow action yields the general nonnegative function
\begin{equation}
    f_{\varphi}(\Delta\varphi).
\end{equation}

The sine form is recovered on a restricted oriented branch when
\begin{equation}
    \boxed{ f_{\varphi}(\Delta\varphi) = \sin(\Delta\varphi), \qquad 0 \leq \Delta\varphi \leq \pi }.
\end{equation}

Within this interval,
\begin{equation}
    \sin(\Delta\varphi) \geq 0.
\end{equation}

Outside this branch, the sine function changes sign and cannot represent a globally positive kinetic coefficient.

A globally defined phase-participation law must therefore be nonnegative, periodic, and compatible with the transformation properties of the underlying topology. Its exact form must be derived from the nonlinear phase-contact action.

Until that derivation is completed, the phase dependence of the hydrodynamic mass should be written as
\begin{equation}
    \boxed{ f_{\varphi}(\Delta\varphi) },
\end{equation}
with the sine relation retained as a branch-level closure.

\subsection{Universal and Mode-Dependent Projections}

The microscopic projection may depend on the defect class and dynamical state:
\begin{equation}
    \iota_V^{(\mathrm{def})} = \iota_V^{(\mathrm{def})} \left( r_{\mathrm{D}}, r_{\mathrm{amp}}, \Delta\varphi, \omega, \text{topology}, \text{boundary conditions} \right).
\end{equation}

For a harmonic mode of angular frequency $\omega$, the conservative inertial projection is determined by the low-frequency real limit
\begin{equation}
    \boxed{ \iota_{V,0}^{(\mathrm{def})} = \lim_{\omega\rightarrow 0} \operatorname{Re} \left[ \iota_V^{(\mathrm{def})}(\omega) \right] }.
\end{equation}

At finite frequency, the response may be dispersive. When dissipative or radiative mechanisms are included, the impedance may acquire an imaginary component.

A universal condensate impedance density exists only if
\begin{equation}
    \iota_{V,0}^{(\mathrm{def})}
\end{equation}
has the same value for all stable defects within the relevant universality class.

In that case, one may define
\begin{equation}
    \boxed{ \rho_{\mathrm{IC}} \equiv \mu_{\mathrm{D}}\iota_{V,0} }.
\end{equation}

If the projection depends on topology or mode structure, then
\begin{equation}
    \rho_{\mathrm{IC}}^{(\mathrm{imp})}
\end{equation}
must remain a mode-resolved effective impedance density rather than being interpreted as a universal intrinsic material density.

\subsection{Primary Added-Inertia Functional}

The factorized form
\begin{equation}
    m_{\mathrm{add}} = \rho_{\mathrm{IC}} V_{\mathrm{eff}} f_{\varphi}
\end{equation}
is valid only when the condensate response separates into:
\begin{equation}
    \text{a universal projected density},
\end{equation}
\begin{equation}
    \text{a geometric or mode participation volume},
\end{equation}
and
\begin{equation}
    \text{a phase-participation factor}.
\end{equation}

The parent relative-flow action does not require such a separation.

For a uniform phase mismatch, the primary surface expression is
\begin{equation}
    \boxed{ m_{\mathrm{add}} = \mu_{\mathrm{D}} \int_{\partial\Omega_{\mathrm{def}}} \left| \mathbf{b}_{\partial} \right|^2 f_{\varphi}(\Delta\varphi) \, \mathrm{d}A }.
\end{equation}

The equivalent volume expression is
\begin{equation}
    \boxed{ m_{\mathrm{add}} = \rho_{\mathrm{IC}}^{(\mathrm{imp})} \int_{\Omega_{\mathrm{IC}}^{(\mathrm{def})}} \left| \boldsymbol{\Upsilon} \right|^2 f_{\varphi}(\Delta\varphi) \, \mathrm{d}V }.
\end{equation}

For a nonuniform phase field, the local quantity
\begin{equation}
    f_{\varphi} \left( \Delta\varphi(\mathbf{x}) \right)
\end{equation}
must remain inside the integrals.

If the phase weighting, condensate flow, and geometric amplitude cannot be separated, the added inertia must remain in this integral form. The factorized hydrodynamic mass law is then only a reduced approximation.

\subsection{Distinction from the Coherent Acoustic Projection}

The coherent FCC contact-cell calculation established
\begin{equation}
    \rho_{\mathrm{eff}}^{(0)} = \mu_{\mathrm{D}}\iota_{\mathrm{tr}},
\end{equation}
where $\iota_{\mathrm{tr}}$ is the surface-to-volume projection associated with a periodic transverse acoustic mode.

The present derivation gives
\begin{equation}
    \rho_{\mathrm{IC}}^{(\mathrm{imp})} = \mu_{\mathrm{D}}\iota_V^{(\mathrm{def})},
\end{equation}
where $\iota_V^{(\mathrm{def})}$ is associated with localized condensate flow relative to the FCC scaffold.

The two projections represent different mechanical operations:
\begin{equation}
    \boxed{ \iota_{\mathrm{tr}} = \text{coherent primitive-cell transverse projection} },
\end{equation}
whereas
\begin{equation}
    \boxed{ \iota_V^{(\mathrm{def})} = \text{localized defect--condensate impedance projection} }.
\end{equation}

Their equality is not implied:
\begin{equation}
    \boxed{ \iota_V^{(\mathrm{def})} \neq \iota_{\mathrm{tr}} }
\end{equation}
unless the corresponding response integrals are shown to coincide.

The coherent acoustic inertia and the localized defect impedance must therefore remain distinct.

\subsection{Consequences for the Existing Calibration}

The assignment
\begin{equation}
    \iota_V = 1 \ \mathrm{m}^{-1}
\end{equation}
must be interpreted as an empirical estimate of the low-frequency defect projection
\begin{equation}
    \iota_{V,0}^{(\mathrm{def})} = \frac{1}{ \ell_{\mathrm{imp}}^{(\mathrm{def})} }.
\end{equation}

The resulting numerical values of
\begin{equation}
    \rho_{\mathrm{IC}}, \qquad K_{\mathrm{D}}, \qquad r_{\mathrm{amp}}, \qquad r_{\mathrm{D}}
\end{equation}
remain conditional on the corresponding assumed penetration length.

A parameter-free determination requires solution of the conservative boundary-value problem for
\begin{equation}
    \boldsymbol{\Upsilon}(\mathbf{x})
\end{equation}
in the actual FCC interstitial geometry, followed by evaluation of
\begin{equation}
    A_{\mathrm{imp}}, \qquad V_{\mathrm{imp}}, \qquad \iota_V^{(\mathrm{def})}, \qquad \mathcal{G}_{\mathrm{imp}}, \qquad f_{\varphi}.
\end{equation}

The present section replaces the dimensional arbitrariness of $\iota_V$ with a definite mechanical functional. It does not assign a numerical value before the corresponding localized flow field has been solved.

\subsection{Constitutive Result}

The localized defect induces a condensate velocity relative to the FCC scaffold:
\begin{equation}
    \boxed{ \mathbf{v}_{\mathrm{rel}} = \dot{q} \boldsymbol{\Upsilon} }.
\end{equation}

On the zero-circulation conservative branch, the response field satisfies
\begin{equation}
    \boxed{ \nabla\cdot\boldsymbol{\Upsilon} = 0, \qquad \nabla\times\boldsymbol{\Upsilon} = \mathbf{0} },
\end{equation}
with boundary condition
\begin{equation}
    \boxed{ \boldsymbol{\Upsilon}\cdot\hat{\mathbf{n}} = \mathbf{b}_{\partial}\cdot\hat{\mathbf{n}} }.
\end{equation}

The mode-weighted interface area and condensate volume are
\begin{equation}
    \boxed{ A_{\mathrm{imp}} = \int_{\partial\Omega_{\mathrm{def}}} \left| \mathbf{b}_{\partial} \right|^2 \, \mathrm{d}A },
\end{equation}
and
\begin{equation}
    \boxed{ V_{\mathrm{imp}} = \int_{\Omega_{\mathrm{IC}}^{(\mathrm{def})}} \left| \boldsymbol{\Upsilon} \right|^2 \, \mathrm{d}V }.
\end{equation}

The microscopic impedance projection is
\begin{equation}
    \boxed{ \iota_V^{(\mathrm{def})} = \frac{ A_{\mathrm{imp}} }{ V_{\mathrm{imp}} } = \frac{1}{ \ell_{\mathrm{imp}}^{(\mathrm{def})} } }.
\end{equation}

The corresponding mode-resolved condensate impedance density is
\begin{equation}
    \boxed{ \rho_{\mathrm{IC}}^{(\mathrm{imp})} = \mu_{\mathrm{D}} \iota_V^{(\mathrm{def})} }.
\end{equation}

The localized added inertia is
\begin{equation}
    \boxed{ m_{\mathrm{add}} = \rho_{\mathrm{IC}}^{(\mathrm{imp})} V_{\mathrm{imp}} f_{\varphi}(\Delta\varphi) = \mu_{\mathrm{D}} A_{\mathrm{imp}} f_{\varphi}(\Delta\varphi) }.
\end{equation}

For a thin spherical defect layer,
\begin{equation}
    V_{\mathrm{sweep}} = 4\pi r_{\mathrm{D}}^2 r_{\mathrm{amp}} + \mathcal{O} \left( r_{\mathrm{D}}r_{\mathrm{amp}}^2 \right),
\end{equation}
and
\begin{equation}
    V_{\mathrm{imp}} = \mathcal{G}_{\mathrm{imp}} V_{\mathrm{sweep}}.
\end{equation}

The reduced mass law is therefore
\begin{equation}
    \boxed{ m_{\mathrm{add}} = K_{\mathrm{D}}^{(\mathrm{def})} r_{\mathrm{amp}} f_{\varphi}(\Delta\varphi) },
\end{equation}
where
\begin{equation}
    \boxed{ K_{\mathrm{D}}^{(\mathrm{def})} = 4\pi r_{\mathrm{D}}^2 \mu_{\mathrm{D}} \iota_V^{(\mathrm{def})} \mathcal{G}_{\mathrm{imp}} }.
\end{equation}

The earlier relations
\begin{equation}
    \rho_{\mathrm{IC}} = \mu_{\mathrm{D}}\iota_V,
\end{equation}
\begin{equation}
    K_{\mathrm{D}} = 4\pi r_{\mathrm{D}}^2\rho_{\mathrm{IC}},
\end{equation}
and
\begin{equation}
    m = \rho_{\mathrm{IC}} V_{\mathrm{eff}} \sin(\Delta\varphi)
\end{equation}
are recovered only under the simultaneous reduced closures
\begin{equation}
    \iota_V^{(\mathrm{def})} = \text{a universal constant},
\end{equation}
\begin{equation}
    \mathcal{G}_{\mathrm{imp}} = 1,
\end{equation}
and
\begin{equation}
    f_{\varphi}(\Delta\varphi) = \sin(\Delta\varphi)
\end{equation}
on the selected nonnegative phase branch.

The fundamental microscopic law is the localized relative-flow added-inertia functional. The projected condensate density, displaced-volume relation, and sine phase factor are reduced constitutive forms whose validity must be established by the solved defect geometry and nonlinear phase-contact dynamics.
